% This chapter presents unpublished joint work with Reid Barton \cite{barton2026}.
Many categories, such as $\Top$ and $\Cat$,
are not locally Cartesian closed,
but still have certain exponentiable objects and morphisms
\cite{niefield1982, giraud1964}.
\begin{definition}[Exponentiable map]\label{def:exp-map}
  Let $\CC$ be a category with pullbacks.
  A morphism $f : Y \to X$ in $\CC$ is \emph{exponentiable}
  when the pullback functor $f^* : \CC / X \to \CC / Y$
  has a right adjoint $\Pi_f : \CC / Y \to \CC / X$.
  The functor $\Pi_f$ is called the \emph{pushforward along $f$}.
\end{definition}
However, for analysing models of type theory,
it is important to have a more fine-grained condition of being exponentiable
relative to a class of maps $\RR$ in a category $\CC$.

For models of MLTT,
such conditions have been considered in \cite{taylor1987} and \cite{joyal2017},
in the form of $\pi$-clans (using the terminology of \cite{joyal2017}).
Uemura introduced categories with representable maps (CwRs) \cite{uemura2023}
for models of MLTT in a lex category (finite limit category).
Exponentiability conditions for $\pi$-clans are in some ways more restrictive than the
theory of exponentiable morphisms --
the example of $\Cat$ and $\Top$ are lost.
At the same time,
it introduces certain closure conditions on pushforwards --
that pushforwards preserve morphisms in the $\pi$-clan,
key to identifying a class of type families in a model.
On the other hand, the above definition of exponentiability (used in CwRs)
captures when a class of maps is exponentiable with respect to \emph{all maps},
which can recover some of examples in $\Cat$ and $\Top$,
but does not subsume the theory of $\pi$-clans
(since not all $\pi$-clans are lex)
This notion of exponentiability says nothing about the
preservation of certain morphisms upon application
of the pushforward functor.

Here, we consider morphisms that are only exponentiable with respect to
a class of maps.
Our setting not only subsumes the
definitions of exponentiability in lex categories and $\pi$-clans,
but also provides a new framework for understanding exponentiability conditions for (op)fibrations in $\Cat$,
open and closed embeddings in $\Top$,
tidy and \'etale maps of topological spaces and toposes,
and Kan fibrations in simplicial sets $\sSet$,
to name but a few examples.
This work will focus on applying this theory to
produce a study of $\Cat$ in terms of \emph{algebraic type theory},
which is the focus of \Cref{sec:cat-as-types}.
% From these considerations,
% we can glean a new kind of dependent type theory
% with models in $\Cat$, $\Top$ and so on.
% This is the focus 

\section{Preclans}

The class of maps with respect to which we consider exponentiability conditions
are required to satisfy certain basic conditions,
which have been considered in various different contexts.
Paul Taylor introduced the notion of \emph{display maps}
for the semantics of Martin-L\"of Type theory,
which could be required to satisfy various extra conditions,
each corresponding to a particular type-theoretic notion
\cite{taylor1987}.
Andr\'e Joyal called a particular combination of those conditions
a \emph{clan} \cite{joyal2017}.
For the sake of having a short name,
we introduce another abbreviation,
consistent with the ``clan'' nomenclature.

\begin{definition}\label{def:preclan}
  A \emph{preclan} $(\CC,\RR)$ consists of a category $\CC$,
  and a class of morphisms $\RR$ in $\CC$
  -- which we call $\RR$-maps -- such that 
  \begin{itemize}
    \item $\RR$ is stable under pullback,
      i.e. pullbacks of $\RR$-maps along all morphisms exist and are $\RR$-maps.
    \item All isomorphisms are in $\RR$.
    \item $\RR$ is closed under composition.
  \end{itemize}
  A \emph{preclan morphism} 
  $F : (\CC, \RR) \to (\CC', \RR')$ is a functor
  $F : \CC \to \CC'$ that
  \begin{itemize}
    \item sends $\RR$-maps to $\RR'$-maps and
    \item preserves pullbacks of $\RR$-maps,
    meaning that for any $\RR$-map $A \to X$
    and any map $Y \to X$,
    the canonical comparison map
    $F (A \times_X Y) \to F A \times_{F X} F Y$
    is an isomorphism.
  \end{itemize}
\end{definition}
Preclans, preclan morphisms,
and natural transformations naturally form a 2-category.

If we also require that $\CC$ has a terminal object
and that maps to terminal objects belong to $\RR$,
then $(\CC,\RR)$ is a \emph{clan}.
A \emph{clan morphism} is a preclan morphism between clans
that also preserves the terminal object.

The full subcategory of $\RR$-objects in a slice
-- objects whose underlying map is an $\RR$-map -- is denoted by
$\RR(X) \subseteq \CC / X$.

Like locally Cartesian closed categories,
preclans are meant to be somewhat type theoretic.
We should therefore expect that each slice
of a preclan is a preclan.
There are in fact two candidates for slicing over an object $X$,
one is $\CC / X$ and the other is $\RR(X)$.

\begin{definition}
  Let $(\CC, \RR)$ be a preclan.
  \begin{itemize}
  \item Given an object $X \in \CC$, the slice category $\CC / X$
    has a preclan structure where the chosen class of maps consists of
    morphisms in $\CC / X$ with an underlying morphism in $\RR$.
    By slight abuse of notation, we denote this preclan by $(\CC/X, \RR)$.
  \item The restriction of the above preclan structure to $\RR(X) \subseteq \CC/X$
    defines a clan $(\RR(X), \RR)$.
  \item
    The category $\wCC$ of presheaves has a preclan structure $(\wCC, \wRR)$
    \cite[Proposition 1.6.5]{joyal2017}.
    Here, $\wRR$ consists of $\RR$-representable natural transformations:
    natural transformations $A \to X$
    such that for any representable object $Y \in \CC$ and map $f : Y \to X$,
    there is an $\RR$-map $A' \to Y$ in $\CC$ representing the pullback
    \[\begin{tikzcd}
        {A'} \pbcorner & A \\
        Y & X
        \arrow[from=1-1, to=1-2]
        \arrow["{\RR\text{-map}}"', from=1-1, to=2-1]
        \arrow[from=1-2, to=2-2]
        \arrow["f"', from=2-1, to=2-2]
      \end{tikzcd}\]
  \end{itemize}
  For an object $X \in \CC$,
  there are inclusions
  \[
    \RR(X) \,\subseteq\, \CC/X \,\subseteq\, \widehat{\CC/X} \simeq \wCC/X
  \]
  These form \emph{preclan embeddings}
  $(\RR(X),\RR) \hookrightarrow \cdots \hookrightarrow (\wCC/X, \wRR)$
  in that they are fully faithful preclan morphisms that reflect $\RR$-maps.
  Furthermore, the Yoneda embedding induces an equivalence $\RR(X) \simeq \wRR(X)$,
  as a representable natural transformation with representable codomain
  also has representable domain.
\end{definition}

\begin{remark}
  For $(\RR(X), \RR)$ to be a preclan, we needed that $\RR$ is closed under composition.
  Pullbacks in $\CC / X$ are computed as pullbacks of the underlying maps in $\CC$.
  Since the inclusion $\RR(X) \to \CC / X$
  reflects pullbacks and $\RR$ is closed under composition,
  the same is also true for $\RR(X)$.
  Moreover, for $(\CC,\RR)$ satisfying only conditions
  $(1)$ and $(2)$ from the definition of preclans,
  $\RR$ is closed under composition if and only if
  $(\RR(X),\RR)$ satisfies $(1)$.
\end{remark}

\begin{definition}\label{def:partial-right-adjoint}
  A \emph{partial right adjoint} of a functor $L : \CC \to \DD$
  consists of two full subcategories
  $\partial\CC \hookrightarrow \CC$
  and $\partial\DD \hookrightarrow \DD$,
  a functor $\partial R : \partial \DD \to \partial \CC$,
  and a hom-set bijection for all
  $X \in \CC$ and $Y \in \partial\DD$
  \[ \DD(L (X), Y) \iso \CC(X,\partial R (Y))\]
  natural in both $X$ and $Y$.
  We depict such a partial right adjoint in the following manner:
  \[\begin{tikzcd}
	\CC & {\partial\CC} \\
	\DD & {\partial \DD}
	\arrow[""{name=0, anchor=center, inner sep=0}, "L"', from=1-1, to=2-1]
	\arrow[hook', from=1-2, to=1-1]
	\arrow[""{name=1, anchor=center, inner sep=0}, "{\partial R}"', from=2-2, to=1-2]
	\arrow[hook', from=2-2, to=2-1]
	\arrow["{\dashv_\partial}"{description}, draw=none, from=0, to=1]
  \end{tikzcd}\]
\end{definition}

For any $\RR$-map $f : Y \to X$ in a preclan $(\CC,\RR)$,
we have a functor $f_! : \CC / Y \to \CC / X$
between the slice categories given by composition with $f$,
with pullback as a partial right adjoint
\[\begin{tikzcd}
{\CC / Y} & {\RR(Y)} \\
{\CC / X} & {\RR(X)}
\arrow[""{name=0, anchor=center, inner sep=0}, "{f_!}"', from=1-1, to=2-1]
\arrow[hook', from=1-2, to=1-1]
\arrow[""{name=1, anchor=center, inner sep=0}, "{f^*}"', from=2-2, to=1-2]
\arrow[hook', from=2-2, to=2-1]
\arrow["{\dashv_\partial}"{description}, draw=none, from=0, to=1]
\end{tikzcd}\]

\begin{proposition}\label{prop:beck-chevalley-left}
  Let $f : Y \to X$ be an $\RR$-map in a preclan $(\CC,\RR)$.
  Then $f_! : \CC / Y \to \CC / X$ restricts to a functor
  $f_! : \RR(Y) \to \RR(X)$.
  Consider the following pullback square.
  \[
  \begin{tikzcd}
    {Y'} & Y \\
    {X'} & X
    \arrow["{s'}", from=1-1, to=1-2]
    \arrow["{f'}"', from=1-1, to=2-1]
    \arrow["\lrcorner"{anchor=center, pos=0.125}, draw=none, from=1-1, to=2-2]
    \arrow["f", from=1-2, to=2-2]
    \arrow["s"', from=2-1, to=2-2]
  \end{tikzcd}
  \]
  The associated Beck--Chevalley transformation
  $(f')_! (s')^* \to s^* f_!$ is a natural isomorphism.
  \[
  \begin{tikzcd}
    \RR(Y') \ar[d, "f'_!"'] \ar[dr, phantom, "\cong"]  &
    \RR(Y) \ar[l, "(s')^*"'] \ar[d, "f_!"] \\
    \RR(X') & \RR(X) \ar[l, "s^*"']
  \end{tikzcd}
  \]
\end{proposition}

\section{The Beck--Chevalley condition}

Variations of the following ``Beck--Chevalley condition''
are typically considered when requiring a right adjoint to pullback.
% In lex categories the ``Beck--Chevalley condition''
% holds whenever the right adjoint to pullback exists,
% making it a theorem, rather than a condition.
% In models of MLTT with $\Pi$-types,
% the ``Beck--Chevalley condition'' is automatically stable under pullback.
% We will see that for us,
% not all right adjoints satisfy the ``Beck--Chevalley condition'',
% and sometimes Beck--Chevalley holds for $f$,
% but not its pullbacks.

\begin{definition} \label{def:beck-chevalley}
  Let $f : Y \to X$ be a morphism in a preclan $(\CC,\RR)$.
  We will say that the \emph{(pushforward) Beck--Chevalley condition}
  holds at $f : Y \to X$ if for any pullback square
   \[ \begin{tikzcd}
        Y' \ar[r, "s'"] \ar[d, "f'"'] \pbcorner & Y \ar[d, "f"] \\
        X' \ar[r, "s"] & X
      \end{tikzcd} \]
    the pullback functor
    $(f')^* : \RR(X') \to \RR(Y')$
    has a right adjoint $f'_* : \RR(Y') \to \RR(X')$,
    and the associated Beck--Chevalley transformation
    $s^* f_* \to (f')_* (s')^*$
    is a natural isomorphism:
    \[
      \begin{tikzcd}
        \RR(Y') \ar[d, "f'_*"'] \ar[dr, phantom, "\cong"]  &
        \RR(Y) \ar[l, "(s')^*"'] \ar[d, "f_*"] \\
        \RR(X') & \RR(X) \ar[l, "s^*"']
      \end{tikzcd}
    \]
\end{definition}

Unlike in locally Cartesian closed categories,
the pushforward Beck--Chevalley condition does not follow
from a Beck--Chevalley isomorphism $(s')_! (f')^* \iso f^* s_!$
(like in \Cref{prop:beck-chevalley-left}) between left adjoints,
because the left adjoint to pullback
$s_! : \RR(X') \to \RR(X)$ may not even exist.
Indeed, we will see in \Cref{ex:topos-etale} that the
pushforward pushforward Beck--Chevalley condition
does not hold for all maps, 
even though we always have \Cref{prop:beck-chevalley-left}.

% We provide an equivalent formulation
% of the Beck--Chevalley condition
% (henceforth we drop the word ``pushforward'').

\begin{proposition} \label{prop:pra-defs}
  Let $f : Y \to X$ be a morphism in a preclan $(\CC,\RR)$.
  % (In particular $f^* : \RR(X) \to \RR(Y)$
  % has a right adjoint $f_* : \RR(Y) \to \RR(X)$.)
  The following conditions are equivalent
  \begin{itemize}
  \item
    The pullback functor
    $f^* : \wCC / X \to \wCC / Y$ between slices in the category of presheaves
    has a partial right adjoint (\Cref{def:partial-right-adjoint})
    $f_* : \RR(Y) \to \RR(X)$.
    \[\begin{tikzcd}
      \wCC / X & {\RR(X)} \\
      \wCC / Y & {\RR(Y)}
      \arrow[""{name=0, anchor=center, inner sep=0}, "f^*"', from=1-1, to=2-1]
      \arrow[hook', from=1-2, to=1-1]
      \arrow[""{name=1, anchor=center, inner sep=0}, "{f_*}"', dashed, from=2-2, to=1-2]
      \arrow[hook', from=2-2, to=2-1]
      \arrow["{\dashv_\partial}"{description}, draw=none, from=0, to=1]
    \end{tikzcd}\]
    
    % This means that we have
    % \[ \Hom_{\CC / Y}(f^* X', A) \iso \Hom_{\CC / X}(X', f_* A) \]
    % natural in $X' \in (\CC / X)^\op$ and $A \in \RR(Y)$.
    % In short, $f_* A$
    % is universal among all objects in the slice $\CC / X$,
    % rather than just the objects in $\RR(X)$.
  \item
    Using the Yoneda embedding $\CC \subseteq \wCC$,
    the right adjoint $\Pi_f : \wCC/Y \to \wCC/X$
    to pullback in presheaves $f^* : \wCC/X \to \wCC/Y$
    sends $\RR(Y) \subseteq \wCC/Y$
    into $\RR(X) \subseteq \wCC/X$.
    \[\begin{tikzcd}
      {\wCC / X} & {\RR(X)} \\
      {\wCC / Y} & {\RR(Y)}
      \arrow[hook', from=1-2, to=1-1]
      \arrow["{\Pi_f}", from=2-1, to=1-1]
      \arrow["{f_*}"', dashed, from=2-2, to=1-2]
      \arrow[hook', from=2-2, to=2-1]
    \end{tikzcd}\]
  \end{itemize}
\end{proposition}
\begin{proof}
  The partial right adjoint exists if and only if
  for any $A \in \RR(Y)$ there exists an object $B \in \RR(X)$
  and a natural isomorphism
  \[ \wCC/Y (f^*(-),A) \iso \wCC/X((-),B)\]
  between functors of type $(\wCC / X)^\op \to \Set$.  
  This holds if and only if
  there exists an object $B \in \RR(X)$ such that $\Pi_f A \iso B$,
  if and only if $\Pi_f$ preserves $\RR$-objects.
\end{proof}

\begin{definition} 
  We say the \emph{partial right adjoint condition}
  holds at $f : Y \to X$ if
  $f$ satisfies the equivalent conditions of
  \Cref{prop:pra-defs}.
\end{definition}

If $\CC$ is lex then the partial right adjoint condition
is equivalent to saying that the pullback
$f^* : \CC / X \to \CC / Y$
in $\CC$ has a partial right adjoint
$f_* : \RR(Y) \to \RR(X)$.
If, additionally, $f$ is exponentiable
(\Cref{def:exp-map})
then the partial right adjoint condition is also equivalent to saying that
the right adjoint $\Pi_f : \CC/Y \to \CC/X$ sends $\RR(Y)$ into $\RR(X)$.
The point is that although the functors
$f^* : \CC / X \to \CC / Y$ and $\Pi_f : \CC/Y \to \CC/X$
may not always exist,
when they do exist,
we can use them instead of their counterparts in $\wCC$.
\[
  \begin{tikzcd}
	{\RR(Y)} & {\CC / Y} & {\wCC / Y} \\
	{\RR(X)} & {\CC / X} & {\wCC / X}
	\arrow[hook, from=1-1, to=1-2]
	\arrow[""{name=0, anchor=center, inner sep=0}, "{f^*}"', shift right=3, from=1-1, to=2-1]
	\arrow[hook, from=1-2, to=1-3]
	\arrow[""{name=1, anchor=center, inner sep=0}, "{f^*}"', shift right=3, dashed, from=1-2, to=2-2]
	\arrow[""{name=2, anchor=center, inner sep=0}, "{f^*}"', shift right=3, from=1-3, to=2-3]
	\arrow[""{name=3, anchor=center, inner sep=0}, "{f_*}"', shift right=3, from=2-1, to=1-1]
	\arrow[hook, from=2-1, to=2-2]
	\arrow[""{name=4, anchor=center, inner sep=0}, "{\Pi_f}"', shift right=3, dashed, from=2-2, to=1-2]
	\arrow[hook, from=2-2, to=2-3]
	\arrow[""{name=5, anchor=center, inner sep=0}, "{\Pi_f}"', shift right=3, from=2-3, to=1-3]
	\arrow["\dashv"{anchor=center}, draw=none, from=0, to=3]
	\arrow["\dashv"{anchor=center}, draw=none, from=1, to=4]
	\arrow["\dashv"{anchor=center}, draw=none, from=2, to=5]
\end{tikzcd}\]

Next, we show that the Beck--Chevalley condition is equivalent
to the partial right adjoint condition.
This was originally observed by Thomas Streicher,
and can be found in \cite[Theorem 1.28]{streicher2012semantics}.
However,
our setting is different to Streicher's in a very subtle way.
Streicher, following Paul Taylor,
requires that $f$ itself is an $\RR$-map, whereas we do not.
We will see that the Beck--Chevalley condition in our setting
is \emph{not necessarily stable under pullback},
unlike in Taylor's setting.
The examples and \Cref{theorem:exponentiable}
will show that throwing in pullback-stability for Beck--Chevalley
actually produces a very natural notion of exponentiability.

\begin{lemma} \label{lem:pullback-global-sections}
  For $s : X' \to X$ and $A \in \RR(X)$, we have
  $\Hom_{\CC/X}(X', A) \cong \Hom_{\RR(X')}(1_{X'}, s^* A)$.
\end{lemma}
\begin{proof}
  Write $A = (A \xrightarrow{p} X)$.
  The left side consists of maps $u : X' \to A$ with $pu = s$,
  while the right side consists of sections of the pullback $s^* p : s^* A \to X'$.
  It makes sense to write $\RR(X')$ on the right hand side
  because both $1_{X'} : X' \to X'$ and $s^* p : s^* A \to X'$ are $\RR$-maps.
\end{proof}

\begin{proposition} \label{prop:beck-chevalley-defs}
  Let $f : Y \to X$ be a morphism in a preclan $(\CC,\RR)$.
  Then the Beck--Chevalley condition holds at $f$ if and only if
  the partial right adjoint condition holds at $f$.
\end{proposition}
\begin{proof}
  Since $\RR(X) \simeq \wRR(X)$,
  the Beck--Chevalley condition holds at $f : Y \to X$
  in $(\CC,\RR)$ if and only if the Beck--Chevalley condition holds at
  $f : Y \to X$ in $(\wCC,\wRR)$.
  Furthermore, since $\wCC$ is lex,
  the partial right adjoint condition holds at $f : Y \to X$
  in $(\CC,\RR)$
  if and only if it holds at $\Pi_f : \wCC / X \to \wCC / Y$
  has a partial right adjoint $f_* : \wRR(Y) \to \wRR(X)$.
  Thus, in the following, we assume that $\CC$ is lex,
  and prove that the Beck--Chevalley condition holds at $f$
  if and only if $f^* : \CC / X \to \CC / Y$
  has a partial right adjoint $f_* : \RR(Y) \to \RR(X)$.

  ($\Rightarrow$)
  We need to establish an isomorphism
  $\Hom_{\CC/X}(X', f_* A) \cong \Hom_{\CC/Y}(f^* X', A)$
  for any $(X' \xrightarrow{s} X) \in \CC/X$ and $(A \to Y) \in \RR(Y)$.
  For $Z \in \CC$, write $\Gamma : \RR(Z) \to \Set$
  for the ``global sections'' functor $\Gamma := \Hom_{\RR(Z)}(1_Z, -)$.
  Let $(Y' \xrightarrow{s'} Y) := f^* (X' \xrightarrow{s} X) \in \CC/Y$,
  with $f' : Y' \to X'$ the induced map,
  so there is a pullback square as in (1).
  Using \Cref{lem:pullback-global-sections},
  \[
    \Hom_{\CC/X}(X', f_* A)
    = \Hom_{\RR(X')}(1_{X'}, s^* f_* A)
    = \Gamma(s^* f_* A)
  \]
  and likewise
  \begin{align*}
    \Hom_{\CC/Y}(f^* X', A)
    & = \Hom_{\RR(Y')}(1_{Y'}, (s')^* A) \\
    & = \Hom_{\RR(Y')}((f')^* 1_{X'}, (s')^* A) \\
    & = \Hom_{\RR(X')}(1_{X'}, f'_* (s')^* A) \\
    & = \Gamma(f'_* (s')^* A)
  \end{align*}
  Condition (1) says that the natural map
  $s^* f_* A \to f'_* (s')^* A$ is an isomorphism,
  so we get an induced isomorphism
  \[
    \Hom_{\CC/X}(X', f_* A)
    = \Gamma(s^* f_* A)
    \cong \Gamma(f'_* (s')^* A)
    = \Hom_{\CC/Y}(f^* X', A)
  \]

  ($\Leftarrow$)
  Consider a pullback square of $f$ as shown below.
  \[
    \begin{tikzcd}
      Y' \ar[r, "s'"] \ar[d, "f'"'] \pbcorner & Y \ar[d, "f"] \\
      X' \ar[r, "s"] & X
    \end{tikzcd}
  \]
  Let $A \in \RR(Y)$.
  We must give a natural isomorphism $s^* f_* A \cong (f')_* (s')^* A \in \RR(X')$.
  We will do this by showing that $s^* f_* A$ represents
  the functor $\Hom_{\CC/Y'}((f')^*(-), (s')^*A)$ on $\CC/X'$.
  Indeed,
  \begin{align*}
    \Hom_{\CC/Y'}((f')^*(-), (s')^*A)
    & \cong \Hom_{\CC/Y}((s')_!(f')^*(-), A) \\
    & \cong \Hom_{\CC/Y}(f^* s_!(-), A) \\
    & \cong \Hom_{\CC/X}(s_!(-), f_* A) \\
    & \cong \Hom_{\CC/X'}({-}, s^* f_* A)
  \end{align*}
  In the second step, we used the two pullbacks lemma
  ($f$ has all pullbacks).
  In the remaining steps, we used the facts that
  $(s')^*$, $f_*$, $s^*$ are partial right adjoints
  to $s'_!$, $f^*$, $s_!$ respectively.
\end{proof}

By the proposition, the Beck--Chevalley condition and the partial right adjoint condition
are equivalent.
However, we will continue to use two different names
because it is often easier
to check or apply one condition over the other.

\begin{remark}
  For the Beck--Chevalley condition,
  we have not needed that $\RR$ is closed under composition,
  since we have not needed the preclan structure $(\RR(X),\RR)$.
\end{remark}

\section{Exponentiability}
\begin{theorem} \label{theorem:exponentiable}
  Let $(\CC,\RR)$ be a preclan and $f : Y \to X$ a morphism in $\CC$.
  The following are equivalent.
  \begin{enumerate}
  \item
    The Beck--Chevalley condition (\labelcref{def:beck-chevalley})
    holds at every pullback of $f$ in the preclan $(\CC,\RR)$
    (whenever such pullbacks exist).
  % \item 
  %   The Beck--Chevalley condition holds at every pullback of
  %   (the image of) $f$ in the preclan ($\wCC,\wRR$).
  \item
    The slice pushforward in presheaves $\Pi_f : (\wCC/Y, \wRR) \to (\wCC/X, \wRR)$
    is a preclan morphism.
  \end{enumerate}
  
  If $\CC$ is lex and $f$ is an exponentiable morphism
  (\Cref{def:exp-map}),
  (2) can be rephrased as saying that
  the slice pushforward $\Pi_f : (\CC/Y, \RR) \to (\CC/X, \RR)$ is a preclan morphism.
\end{theorem}

Before proving the theorem,
we provide our main definitions,
as well as some lemmas and remarks.

\begin{definition} \label{def:exponentiable}
  We say $f$ is \emph{$\RR$-exponentiable}
  if it satisfies the equivalent conditions in \Cref{theorem:exponentiable}.
  We write $\Exp_\RR$ for the class of $\RR$-exponentiable maps in $\CC$.
\end{definition}

By condition (1), $\RR$-exponentiable morphisms
are stable under pullback.
Condition (2) says that $\Pi_f : (\wCC/Y, \wRR) \to (\wCC/X, \wRR)$
not only sends $\RR(Y) \subseteq \wCC / Y$ into $\RR(X) \subseteq \wCC / X$,
as in the statement of \Cref{prop:beck-chevalley-defs},
but also preserves all $\RR$-maps in
$\RR(Y)$ and $\CC / Y$, which are subcategories of $\wCC/Y$.
Hence, assuming condition (2),
we have that the pushforward $f_* : (\RR(Y), \RR) \to (\RR(X), \RR)$
is a preclan morphism,
making it a partial right adjoint 
to pullback $f^* : (\wCC/X, \RR) \to (\wCC /Y, \RR)$
in the $2$-category of preclans.
  \[\begin{tikzcd}
    (\wCC / X, \wRR) & {(\RR(X), \RR)} \\
    (\wCC / Y, \wRR) & {(\RR(Y), \RR)}
    \arrow[""{name=0, anchor=center, inner sep=0}, "f^*"', from=1-1, to=2-1]
    \arrow[hook', from=1-2, to=1-1]
    \arrow[""{name=1, anchor=center, inner sep=0}, "{f_*}"', from=2-2, to=1-2]
    \arrow[hook', from=2-2, to=2-1]
    \arrow["{\dashv_\partial}"{description}, draw=none, from=0, to=1]
  \end{tikzcd}\]

As we explain in \Cref{sec:example},
condition~(1) is closely related to the
various notions of a map being ``proper''.
Anel and Weinberger \cite{anel2024} consider such conditions
in a different setting where the ``coefficients'' may not be given by
a full subcategory $\RR(-)$ of the slice $\CC/(-)$
(and when $\CC$ is lex, though this is minor).
In the language of comprehension categories \cite{jacobs1993},
their $\RR$ is a comprehension category over $\CC$
whereas ours is a \emph{full} one.
When both settings coincide,
their ``proper maps'' are our $\RR$-exponentiable maps,
stated via a pullback stable Beck-Chevalley condition.
The equivalence (1) $\Leftrightarrow$ (3) in \Cref{prop:beck-chevalley-defs} 
says that when both settings coincide,
their ``strict proper maps'' are the same as ``proper maps''.
% It seems that at least some of this theory works without the assumption
% of being fully faithful;
% we leave such considerations to future work.

\begin{definition}\label{def:double-preclan}
  Let $(\CC,\RR)$ be a preclan
  and $\HH$ a class of maps in $\CC$.
  We say that $\HH$ is $\RR$-exponentiable when all the
  maps in $\HH$ are $\RR$-exponentiable.

  If $(\CC,\HH)$ is also a preclan, $\HH$ is $\RR$-exponentiable
  and $\RR$ is $\HH$-exponentiable,
  then we will say that the combination
  $(\CC,\RR,\HH)$ is a $\tau$-\emph{preclan}.
\end{definition}

Naturally, if both $(\CC,\RR)$ and $(\CC,\HH)$
are clans then we would call the
combination a $\tau$-\emph{clan}.
Examples of $\tau$-preclans can be found in \Cref{ex:top-opens} and
\Cref{ex:categories-and-fibrations}.

\begin{definition}\label{def:pi-preclan}
  We call $(\CC, \RR)$ a $\pi$-\emph{preclan}
  if all $\RR$-maps are $\RR$-exponentiable.
\end{definition}

This is close to Paul Taylor's definition of a
display map category with dependent products \cite{taylor1987},
and recovers what Andr\'e Joyal calls a $\pi$-clan \cite{joyal2017}
when $(\CC,\RR)$ is a clan.

\begin{remark}[Counterexamples]\label{rmk:counterexample} 
  Exponentiable maps in a lex category
  need not be $\RR$-exponentiable
  (see \Cref{ex:categories-and-fibrations}).
  Conversely, $\RR$-exponentiable maps need not be exponentiable
  with respect to all maps, either.
  Consider the $\pi$-clan $(\CC,\RR)$, where 
  $\CC$ is the category of contexts
  (the syntactic category)
  for Martin-L\"of Type Theory with
  $\Unit$-types, $\Sigma$-types, and $\Pi$-types
  and $\RR$ is the class of type families.
  In fact, in this example, the category $\CC$ is not even lex,
  so that the class of all maps does not even produce a clan.
  
  One might wonder why the definition of an $\RR$-exponentiable morphism
  requires the Beck--Chevalley condition for all pullbacks of $f$,
  rather than just for $f$ itself.
  % Experience with examples and applications shows that
  % forcing the class of $\RR$-exponentiable maps to be pullback-stable
  % leads to a more useful notion.
  We have the following implications:
  \[
    \begin{tikzcd}
      \text{$f : Y \to X$ is $\RR$-exponentiable}
      \ar[d, Rightarrow, "\text{\,(\textit{i})}"] \\
      {\begin{array}{c}
        \text{$f$ satisfies Beck--Chevalley and $f_* : (\RR(Y), \RR) \to (\RR(X), \RR)$ is a preclan morphism}
      \end{array}}
      \ar[d, Rightarrow, "\text{\,(\textit{ii})}"] \\
      \text{$f$ satisfies Beck--Chevalley}
      \ar[d, Rightarrow, "\text{\,(\textit{iii})}"] \\
      \text{$f^* : \RR(X) \to \RR(Y)$ has a right adjoint $f_* : \RR(Y) \to \RR(X)$}
    \end{tikzcd}
  \]
  In general, none of these implications is reversible.
  We provide counterexamples for (\textit{i}), (\textit{ii}) and (\textit{iii})
  respectively in \Cref{ex:top-opens}, \Cref{ex:top-open-maps}
  and \Cref{ex:topos-etale}.
\end{remark}

We now turn to the proof of \Cref{theorem:exponentiable}.
In the proof,
we will need to analyse the functorial action of $\Pi_f : \wCC/Y \to \wCC/X$ on morphisms.
To do so, we will reintroduce \cite[Lemmas 2.3.1 and 2.3.2]{joyal2017}
in the following.
First, recall that for any functor
$F : \CC \to \DD$ and object $X \in \CC$,
there is an induced functor between slices, denoted by
\[ F / X : \CC / X \to \DD / F X\]
which acts on objects and morphisms on the slice by
applying $F$ to the underlying morphisms in $\CC$.

\begin{lemma}[{\cite[Lemma~2.3.1 + dual]{joyal2017}}] \label{lem:adjoints-of-slice-functors}
  Suppose we have an adjunction $L \dashv R : \CC \to \DD$
  and an object $X \in \CC$.
  We denote the component of the counit at $X$ as
  $\epsilon : L R X \to X$.
  Then there is an adjunction $\epsilon_! \circ L / R X \dashv R / X$,
  depicted below.

  On the other hand, suppose we have an object $Y \in \DD$
  and the component $\eta : Y \to R L Y$ of the unit at $Y$
  has pullbacks along all maps,
  so that we have a functor $\eta^* : \DD / R L Y \to \DD / Y$.
  Then there is a dual adjunction $L / Y \dashv \eta^* \circ R / L Y$.
  \[ \begin{tikzcd}
    \CC &&& {\CC / X} && {\CC / L Y} \\
    && {\CC / L R X} &&&& {\DD / R L Y} \\
    \DD &&& {\DD / R X} && {\DD / Y}
    \arrow[""{name=0, anchor=center, inner sep=0}, "R", shift left=3, from=1-1, to=3-1]
    \arrow[""{name=1, anchor=center, inner sep=0}, "{{R / X}}", from=1-4, to=3-4]
    \arrow["{R / L Y}", from=1-6, to=2-7]
    \arrow["{{\epsilon_!}}", from=2-3, to=1-4]
    \arrow["{\eta^*}", from=2-7, to=3-6]
    \arrow[""{name=2, anchor=center, inner sep=0}, "L", shift left=3, from=3-1, to=1-1]
    \arrow["{{L / R X}}", from=3-4, to=2-3]
    \arrow[""{name=3, anchor=center, inner sep=0}, "{L / Y}", from=3-6, to=1-6]
    \arrow["\dashv"{description}, draw=none, from=2-3, to=1]
    \arrow["\dashv"{description}, draw=none, from=2, to=0]
    \arrow["\dashv"{description}, draw=none, from=3, to=2-7]
  \end{tikzcd} \]
\end{lemma}
\begin{proof}
  For the first adjunction,
  let $(a : A \to R X) \in \DD / R X$ and $(b : B \to X) \in \CC / X$.
  A map $t$ in $\CC / X (\epsilon_! \circ L / R X (A), B)$
  corresponds to a diagram in $\CC$
  \[ \begin{tikzcd}
    {L A} & B \\
    {L R X} & X
    \arrow["t", dashed, from=1-1, to=1-2]
    \arrow["{L a}"', from=1-1, to=2-1]
    \arrow["b", from=1-2, to=2-2]
    \arrow["\epsilon"', from=2-1, to=2-2]
  \end{tikzcd} \]
  which transposes along the adjunction $L \dashv R$
  to a diagram in $\DD$
  \[ \begin{tikzcd}
    A & RB \\
    {R X} & {R X}
    \arrow["{\tilde{t}}", dashed, from=1-1, to=1-2]
    \arrow["a"', from=1-1, to=2-1]
    \arrow["Rb", from=1-2, to=2-2]
    \arrow[equals, from=2-1, to=2-2]
  \end{tikzcd} \]
  This is precisely a morphism
  $\tilde{t}$ in $\DD / R X (A, R / Z (B))$.
  One can verify that this bijection is natural in both $A$ and $B$.
  
  For the second adjunction,
  let $A \in \DD / Y$ and $B \in \CC / L Y$.
  A map $t$ in $\CC / L Y (L / Y (A), B)$
  corresponds to a diagram in $\CC$,
  depicted below with its transpose along $L \dashv R$ in $\DD$.
  \[ \begin{tikzcd}
    {L A} & B && A & RB \\
    {L Y} & {L Y} && Y & {R L Y}
    \arrow["t", dashed, from=1-1, to=1-2]
    \arrow["{L a}"', from=1-1, to=2-1]
    \arrow["b", from=1-2, to=2-2]
    \arrow["{\tilde{t}}", dashed, from=1-4, to=1-5]
    \arrow["a"', from=1-4, to=2-4]
    \arrow["Rb", from=1-5, to=2-5]
    \arrow[equals, from=2-1, to=2-2]
    \arrow["\eta"', from=2-4, to=2-5]
  \end{tikzcd} \]
  The transpose is a morphism in
  \[\DD / R L Y (\eta_! (A), R / L Y (B)) \cong \DD / Y (A, \eta^* \circ R / L Y)\]
  This provides the required natural bijection.
\end{proof}

It is well known that
exponentiable morphisms are stable under pullback 
(see \cite[Section 2.1]{weber2015} for example,
which uses Dubuc's adjoint triangle theorem \cite{dubuc2006}).
In the following,
we construct the pushforward functor $\Pi_{f'}$ associated
with a pullback $f'$ of an exponentiable morphism~$f$.

\begin{lemma} \label{lem:pushforward-along-pullback}
  Consider an exponentiable morphism $f : Y \to X$
  in a lex category $\CC$.
  Consider a morphism $s : X' \to X$ and its pullback along $f$.
  \[\begin{tikzcd}
    {Y'} & {X'} & {\Pi_f Y'} \\
    Y & X
    \arrow["{{{f'}}}", from=1-1, to=1-2]
    \arrow["{{{s'}}}"', from=1-1, to=2-1]
    \arrow["\lrcorner"{anchor=center, pos=0.125}, draw=none, from=1-1, to=2-2]
    \arrow["{{{\eta}}}", from=1-2, to=1-3]
    \arrow["s"', from=1-2, to=2-2]
    \arrow["{{{\Pi_f s'}}}", from=1-3, to=2-2]
    \arrow["f"', from=2-1, to=2-2]
  \end{tikzcd}\]
  where $\eta : X' \to \Pi_f Y'$ denotes the unit of the adjunction
  $f^* \dashv \Pi_f$ at $(s : X' \to X)$.
  Then the composition $\Pi_{f'} := \eta^* \circ \Pi_f / Y'$
  is right adjoint to the pullback functor
  $(f')^* : \CC / X' \to \CC / Y'$.
  \[ \begin{tikzcd}
    {\CC / Y'} & {\CC / \Pi_f Y'} & {\CC / X'}
    \arrow["{\Pi_f / Y'}", from=1-1, to=1-2]
    \arrow["{\Pi_{f'}}"', bend right, from=1-1, to=1-3]
    \arrow["{\eta^*}", from=1-2, to=1-3]
  \end{tikzcd} \]
\end{lemma}
\begin{proof}
  The left adjoint of $\eta^* \circ \Pi_f / Y'$
  is $f^* / X' : \CC / X' \to \CC / Y'$,
  by \Cref{lem:adjoints-of-slice-functors}.
  Clearly $(f')^* \cong f^* / X'$, as required. 
\end{proof}

Now suppose that we also have $g : Z \to Y$.
Then the pushforward $\Pi_f$ induces a functor
$\Pi_f / Z : \CC / Z \to \CC / \Pi_f Z$ on slices,
which we can calculate using the following lemma.

\begin{lemma}[{\cite[Lemma~2.3.2]{joyal2017}}] \label{lem:pullback-morphism-action}
  Suppose $f : Y \to X$ is an exponentiable morphism
  in a lex category $\CC$,
  and $g : Z \to Y$.
  Consider the diagram
  \[\begin{tikzcd}
    Z & {Y \times_X \Pi_f Z} & {\Pi_f Z} \\
    & Y & X
    \arrow["g"', from=1-1, to=2-2]
    \arrow["{\epsilon}"', from=1-2, to=1-1]
    \arrow["{f'}", from=1-2, to=1-3]
    \arrow[from=1-2, to=2-2]
    \arrow["\lrcorner"{anchor=center, pos=0.125}, draw=none, from=1-2, to=2-3]
    \arrow["{\Pi_f g}", from=1-3, to=2-3]
    \arrow["f"', from=2-2, to=2-3]
  \end{tikzcd}\]
  where $\epsilon : Y \times_X \Pi_f Z \to Z$
  is the counit of the adjunction $f^* \dashv \Pi_f$ at $(g : Z \to Y)$.
  Then there is a canonical natural isomorphism
  $\Pi_f/Z \iso \Pi_{f'} \circ \epsilon^*$.
  \begin{equation} \label{diag:pullback-morphism-action}
    \begin{tikzcd}
      {\CC / Z} & {\CC / \big( Y \times_X \Pi_f Z \big)} & {\CC / \Pi_f Z}
      \arrow["{{\epsilon^*}}", from=1-1, to=1-2]
      \arrow["{{\Pi_{f'}}}", from=1-2, to=1-3]
      \arrow["{{\Pi_f/Z}}"', from=1-1, to=1-3, bend right]
    \end{tikzcd}
  \end{equation}
\end{lemma}
\begin{proof}
  It suffices to construct an isomorphism between their left adjoints.
  The left adjoint of $\Pi_f / Z$ is computed as
  $\epsilon_! \circ f^* / \Pi_f Z$ by \Cref{lem:adjoints-of-slice-functors}.
  By definition,
  \[f^* / \Pi_f Z \iso (f')^* : \CC / \Pi_f Z \to \CC / \big( Y \times_X \Pi_f Z \big)\]
  By composing adjunctions $\epsilon_! \dashv \epsilon^*$ and $(f')^* \dashv \Pi_{f'}$
  we see that $\epsilon_! \circ (f')^* \dashv \Pi_{f'} \circ \epsilon^*$.
\end{proof}

It follows from \Cref{lem:pullback-morphism-action} that we have a canonical isomorphism
\begin{equation} \label{eq:type-theoretic-axiom-of-choice}
  f_* \circ A_! \iso (f_* A)_! \circ f'_* \circ \epsilon^* 
\end{equation}
This is called the \emph{distributivity law} in \cite{gambino2013}
(for $\Pi$-types and $\Sigma$-types).
It could also be called the \emph{type theoretic axiom of choice}.
In an informal type theoretic syntax, one could render it as follows.
\[
  \prod_{x : X} \sum_{y : Y(x)} A(y) \iso \sum_{f : \Pi_{X} Y} \prod_{x : X} A(f x)
\]

\begin{lemma}\label{lem:BC-iff-BC-in-presheaves}
  Let $(\CC,\RR)$ be a preclan and $f : Y \to X$ a morphism in $\CC$.
  The following are equivalent.
  \begin{itemize}
  \item
    The Beck--Chevalley condition holds at every pullback of
    $f$ in the preclan $(\CC,\RR)$
    (whenever such pullbacks exist).
  \item 
    The Beck--Chevalley condition holds at every pullback of
    (the image of) $f$ in the preclan ($\wCC,\wRR$).
  \end{itemize}
\end{lemma}
\begin{proof}
  ($\Rightarrow$)
  Suppose the partial right adjoint condition
  holds at every pullback of $f$ in $(\CC,\RR)$.
  Consider a pullback of $f$ in $\wCC$
  \[
    \begin{tikzcd}
      Y' \ar[r, "t"] \ar[d, "f'"'] \pbcorner & Y \ar[d, "f"] \\
      X' \ar[r, "s"] & X \rlap{\, .}
    \end{tikzcd}
  \]
  Note that $X'$ and $Y'$ in $\wCC$ may not be representable.
  Since $\wCC$ is locally Cartesian closed,
  it suffices to show that $\Pi_{f'} : \wCC / Y' \to \wCC / X'$
  sends $\wRR(Y')$ into $\wRR(X')$
  (this is a simplification of condition (3) in \Cref{prop:beck-chevalley-defs}).
  Let $q : A \to Y'$ be in $\wRR(Y')$
  and write $r := \Pi_{f'}(q)$.
  To check that $r$ belongs to $\wRR$,
  let $u : \widetilde X \to X'$ be a map with $\widetilde X \in \CC$;
  we must check that $u^*(r)$ is a morphism of $\CC$ that belongs to $\RR$.

  \[
    \begin{tikzcd}[row sep=small, column sep=small]
      %JH: with =small it looked like \tilde{f} is a label for another arrow
      %RB: I decided we didn't really need those "back" arrows anyways
      u^* A \ar[rd, "u^*(q)"'] & & A \ar[rd, "q"'] \\
      & \widetilde Y \ar[rr] \ar[dd, "\tilde f"'] \pbcorner & &
      Y' \ar[rr, "t"] \ar[dd, "f'"'] \pbcorner & & Y \ar[dd, "f"] \\
      u^*(\Pi_{f'} A) \ar[rd, "u^*(r)"'] & &
      \Pi_{f'} A \ar[rd, "r"'] & & \hphantom{B} \\
      & \widetilde X \ar[rr, "u"] & &
      X' \ar[rr, "s"] & & X
    \end{tikzcd}
  \]

  As shown above, write $\widetilde f : \widetilde Y \to \widetilde X$
  for the pullback $u^*(f')$.
  Since $u^* : \wCC / X' \to \wCC / \widetilde X$ preserves pushforwards,
  the morphism $u^*(r)$ is isomorphic to $\Pi_{\tilde f}(u^*(q))$.
  Now as $q \in \wRR$ and $\widetilde Y \in \CC$,
  the map $u^*(q)$ belongs to $\RR$.
  The composition $su : \widetilde X \to X$ is a morphism of $\CC$,
  so the morphism $\widetilde f : \widetilde Y \to \widetilde X$
  is a pullback of $f$ in $\CC$ (in particular, $\widetilde Y \in \CC$).
  By the hypothesis on $f$, then,
  the object $u^*(r) \cong \Pi_{\tilde f}(u^*(q))$ of $\wCC/\widetilde X$
  belongs to $\RR(\widetilde X)$.
  Thus, every pullback of $r$ to a representable $\widetilde X$
  produces a morphism in $\RR$, so $r \in \wRR$.
  
  ($\Leftarrow$)
  Now suppose the partial right adjoint condition
  holds at every pullback of $f$ in $(\wCC,\wRR)$.
  Consider a pullback
  \[
    \begin{tikzcd}
      Y' \ar[r, "t"] \ar[d, "f'"'] \pbcorner & Y \ar[d, "f"] \\
      X' \ar[r, "s"] & X
    \end{tikzcd}
  \]
  of $f$ in $\CC$.
  We must show that $\Pi_{f'} : \wCC / Y' \to \wCC / X'$
  sends $\RR(Y')$ into $\RR(X')$
  (condition (3) in \Cref{prop:beck-chevalley-defs}).
  Since the image of the given pullback
  is still a pullback in $\wCC$,
  our assumption says that $\Pi_{f'} : \wCC / Y' \to \wCC / X'$
  sends $\wRR(Y')$ into $\wRR(X')$.
  We conclude by recalling that
  $\RR(Z) \simeq \wRR(Z)$ for any $Z$ in $\CC$.
\end{proof}

\begin{proposition}\label{prop:lex-exponentiable}
  Let $(\CC,\RR)$ be a preclan such that
  $\CC$ is lex and $f : Y \to X$
  an exponentiable morphism in $\CC$.
  The following are equivalent
  \begin{itemize}
  \item
    The Beck--Chevalley condition holds at every pullback of
    $f$ in the preclan $(\CC,\RR)$.
  \item
    The slice pushforward in presheaves $\Pi_f : (\CC/Y, \RR) \to (\CC/X, \RR)$
    is a preclan morphism.
  \end{itemize}
\end{proposition}
\begin{proof}
  % We use
  % the partial right adjoint condition
  % rather than the Beck--Chevalley condition.

  ($\Rightarrow$)
  Suppose $\CC$ is lex, $f : Y \to X$ is exponentiable, and
  the partial right adjoint condition holds at every pullback of $f$.
  We must show that $\Pi_f : \CC / Y \to \CC / X$ is a preclan morphism.
  Note that $\Pi_f : \CC / Y \to \CC / X$ is a right adjoint,
  thus it preserves all pullbacks.
  It remains to show that for any $\RR$-map $h : A \to Z$ over $Y$,
  \[\begin{tikzcd}
    A & Z & Y & X
    \arrow["h", from=1-1, to=1-2]
    \arrow["g", from=1-2, to=1-3]
    \arrow["f", from=1-3, to=1-4]
  \end{tikzcd}\]
  the pushforward $\Pi_f h : \Pi_f A \to \Pi_f Z$
  is also an $\RR$-map.

  By \Cref{lem:pullback-morphism-action},
  we can analyse $\Pi_f h$ by viewing $h : A \to Z$ as an object in $\CC / Z$.
  In our case we can further view $h : A \to Z$ as an object in $\RR(Z)$
  and use that $\Pi_{f'}$ restricts to $f'_*$,
  extending \cref{diag:pullback-morphism-action}.
  \[\begin{tikzcd}
    {\RR(Z)} & {\RR(Y \times_X \Pi_f Z)} & {\RR(\Pi_f Z)} \\
    {\CC / Z} & {\CC / \big( Y \times_X \Pi_f Z \big)} & {\CC / \Pi_f Z}
    \arrow["{\epsilon^*}", from=1-1, to=1-2]
    \arrow[from=1-1, to=2-1]
    \arrow["{f'_*}", from=1-2, to=1-3]
    \arrow[from=1-2, to=2-2]
    \arrow[from=1-3, to=2-3]
    \arrow["{\epsilon^*}", from=2-1, to=2-2]
    \arrow["{\Pi_{f'}}", from=2-2, to=2-3]
    \arrow["{\Pi_f/Z}", from=2-1, to=2-3, bend right]
  \end{tikzcd}\]
  Hence $\Pi_f h$ can be viewed as an object in $\RR(\Pi_f Z)$,
  which is an $\RR$-map in $\CC$:
  \[
    \Pi_f h
    \,\iso\, \Pi_{f'} \, \epsilon^* \, h
    \,\iso\, f'_* \, \epsilon^* \, h \,.
  \]

  ($\Leftarrow$)
  Suppose $\CC$ is lex, $f : Y \to X$ is exponentiable,
  and $\Pi_f : (\CC/Y,\RR) \to (\CC/X,\RR)$ is a preclan morphism.
  Let $f' : Y' \to X'$ be a pullback of $f$
  \[
    \begin{tikzcd}
      Y' \ar[r, "s'"] \ar[d, "f'"'] \pbcorner & Y \ar[d, "f"] \\
      X' \ar[r, "s"] & X
    \end{tikzcd}
  \]
  By \Cref{lem:pushforward-along-pullback} there is
  a pullback functor
  $\Pi_{f'} = \eta^* \circ \Pi_f / Y' : \CC / Y' \to \CC / X'$.
  Since $\Pi_f$ is a preclan morphism,
  $\Pi_f / Y'$ sends $\RR(X')$ into $\RR(\Pi_f Y')$.
  Pullbacks always preserve $\RR$-maps,
  so $\eta^*$ sends $\RR(\Pi_f Y')$ into $\RR(Y')$.
  Hence $\Pi_{f'}$ sends $\RR(Y')$ into $\RR(X')$
  and so the partial right adjoint condition holds at $f'$.
\end{proof}

We can now prove \Cref{theorem:exponentiable},
by combining the previous propositions.

\begin{proof}[Proof of \Cref{theorem:exponentiable}.]
  Let $(\CC,\RR)$ be a preclan and $f : Y \to X$ a morphism in $\CC$.
  The Beck--Chevalley condition holds at every pullback of
  $f$ in the preclan $(\CC,\RR)$
  if and only if 
  it holds at every pullback of $f$ in $(\wCC,\wRR)$,
  by \Cref{lem:BC-iff-BC-in-presheaves}.
  The latter holds if and only if
  the slice pushforward in presheaves $\Pi_f : (\wCC/Y, \wRR) \to (\wCC/X, \wRR)$
  is a preclan morphism,
  by \Cref{prop:lex-exponentiable}
  and the fact that $\wCC$ is locally Cartesian closed.
\end{proof}

The usual definition of a $\pi$-(pre)clan follows.

\begin{corollary} \label{cor:pi-preclan}
  Let $(\CC, \RR)$ be a preclan.
  The following are equivalent
  \begin{enumerate}
  \item
    All $\RR$-maps are $\RR$-exponentiable,
    i.e. $(\CC, \RR)$ is a $\pi$-preclan.
  \item
    The Beck--Chevalley condition holds at all $\RR$-maps.
  \end{enumerate}
  In this case, $\Pi_f : (\wCC / X, \RR) \to (\wCC / Y, \RR)$
  is a preclan morphism for every $\RR$-map $f$,
  and so is $\Pi_f : (\CC / X, \RR) \to (\CC / Y, \RR)$
  when $\CC$ is lex and $f$ is exponentiable.
\end{corollary}
\begin{proof}
  Since $\RR$-maps are stable under pullback,
  a fact about all morphisms in $\RR$ is true if and only if
  it is true for all pullbacks of morphisms in $\RR$.
\end{proof}

\begin{proposition} \label{prop:exp-clan}
  If $(\CC,\RR)$ is a preclan
  then $(\CC,\Exp_\RR)$ is also a preclan.
\end{proposition}
\begin{proof}
  Pullbacks of $\RR$-exponentiable morphisms exist and
  are still $\RR$-exponentiable by considering condition (1)
  from the equivalent definitions of $\RR$-exponentiability.
  If $f$ is an isomorphism, then the Beck--Chevalley condition holds trivially
  at $f$ as well as its pullbacks, which are also isomorphisms.
  If $g : Z \to Y$ and $f : Y \to X$ are both $\RR$-exponentiable,
  then by condition (3)
  \[\Pi_f \circ \Pi_g \iso \Pi_{f \circ g} : (\wCC/Z, \wRR) \to (\wCC/X, \wRR)\]
  we see that $f \circ g$ is also $\RR$-exponentiable.
\end{proof}

The remainder of this section provides various lemmas
about $\RR$-exponentiability that will be useful in the coming chapters.
The first lemma \Cref{lem:stable-class-of-exponentiable-maps}
makes it easier to check that a class of maps $\HH$ is $\RR$-exponentiable,
using a more elementary condition.

\begin{lemma}\label{lem:stable-class-of-exponentiable-maps}
  Suppose $\HH$ is a class of maps that is stable under pullback,
  and for every $f : Y \to X$ in $\HH$,
  and every $\RR$-map $g : Z \to Y$, 
  there is an $\RR$-map $p : P \to X$ satisfying the partial right adjoint
  universal property of the pushforward:
  \[ \CC / Y (f^* h, g) \iso \CC / X (h,p)\]
  naturally for all $h$ in $\CC / X$.
  Then $\HH$ is $\RR$-exponentiable.
  \[\HH \subseteq \Exp_\RR\]
\end{lemma}
\begin{proof}
  Let $f : Y \to X$ be in $\HH$.
  By assumption on $\HH$, $f$ has all pullbacks.
  The elementary condition is sufficient for constructing a
  partial right adjoint $f_* : \RR(Y) \to \RR(X)$
  for the pullback functor $f^* : \CC / X \to \CC / Y$.
  \[\begin{tikzcd}
    \CC / X & {\RR(X)} \\
    \CC / Y & {\RR(Y)}
    \arrow[""{name=0, anchor=center, inner sep=0}, "f^*"', from=1-1, to=2-1]
    \arrow[hook', from=1-2, to=1-1]
    \arrow[""{name=1, anchor=center, inner sep=0}, "{f_*}"', from=2-2, to=1-2]
    \arrow[hook', from=2-2, to=2-1]
    \arrow["{\dashv_\partial}"{description}, draw=none, from=0, to=1]
  \end{tikzcd}\]
  Since $\HH$ is stable under pullback,
  the existence of the partial right adjoint also holds for all pullbacks of $f$.
\end{proof}

Like in \cite[Lemma 2.4.13]{joyal2017},
we also provide versions of 
\Cref{lem:pushforward-along-pullback}
and \Cref{lem:pullback-morphism-action}
for preclan (rather than a lex category).
These lemmas aid in various computations
involving the pushforward,
which we make use of for the theory of polynomial functors in \Cref{sec:polynomials}.

For a preclan morphism $F : (\CC, \RR) \to (\CC',\RR')$
and $X \in \CC$,
the action of $F$ on morphisms induces a functor
\[F / X : \RR(X) \to \RR'(F X)\]
Then for an $\RR$-map $f : X \to Y$
we have $(F f)_! \circ F / X \iso F / Y \circ f_!$.
\[
  \begin{tikzcd}
    {\RR(X)} & {\RR'(F X)} \\
    {\RR(Y)} & {\RR'(F Y)}
    \arrow["{F / X}", from=1-1, to=1-2]
    \arrow["{f_!}"', from=1-1, to=2-1]
    \arrow["{(F f)_!}", from=1-2, to=2-2]
    \arrow["F / Y"', from=2-1, to=2-2]
  \end{tikzcd}
\]

% In particular for 
% \begin{definition}
%   Let $c,d \in \CC$ and $F : (\RR(c), \RR) \to (\RR(d), \RR)$
%   be a preclan morphism.
%   For an object $X$ in $\RR(c)$,
%   the action of $F$ on morphisms induces a functor
%   \[F / X : \RR(X) \to \RR(F X)\]
%   Writing $X_! : \RR(X) \to \RR(c)$ for the functor composing
%   the map $X \to c$,
%   it follows that the following square commutes.
%   \[
%   \begin{tikzcd}
%     {\RR(X)} & {\RR(F X)} \\
%     {\RR(c)} & {\RR(d)}
%     \arrow["{F / X}", from=1-1, to=1-2]
%     \arrow["{X_!}"', from=1-1, to=2-1]
%     \arrow["{(F X)_!}", from=1-2, to=2-2]
%     \arrow["F"', from=2-1, to=2-2]
%   \end{tikzcd}
%   \]
% \end{definition}

\begin{lemma}
  Let $f : Y \to X$ be $\RR$-exponentiable,
  and let $A$ be an object in $\RR(X)$,
  which we pull back along $f$.
  Consider the diagram
  \[\begin{tikzcd}
    {f^* A} & {A} & {f_* f^* A} \\
    Y & X
    \arrow["{{{f'}}}", from=1-1, to=1-2]
    \arrow[from=1-1, to=2-1]
    \arrow["\lrcorner"{anchor=center, pos=0.125}, draw=none, from=1-1, to=2-2]
    \arrow["{{{\eta}}}", from=1-2, to=1-3]
    \arrow[from=1-2, to=2-2]
    \arrow["{{{f_* s'}}}", from=1-3, to=2-2]
    \arrow["f"', from=2-1, to=2-2]
  \end{tikzcd}\]
  where $\eta : A \to f_* f^* A$ denotes the unit of the adjunction
  $f^* \dashv f_*$ at $A$.
  Then there is a canonical isomorphism 
  ${f'_*} \iso \eta^* \circ f_* / f^* A$.
  \[ \begin{tikzcd}
    {\RR(f^* A)} & {\RR(f_* f^* A)} & {\RR(A)}
    \arrow["{f_* / f^* A}", from=1-1, to=1-2]
    \arrow["{{f'_*}}"', bend right, from=1-1, to=1-3]
    \arrow["{\eta^*}", from=1-2, to=1-3]
  \end{tikzcd} \]
\end{lemma}
\begin{proof}
  Embedding everything into presheaves $\RR(A) \to \wCC / A$,
  and noting that all functors commute with the embeddings,
  the result follows from \cite[Lemma 1.13]{barton2026}.
\end{proof}

\begin{lemma}\label{lem:distributivity}
  Let $f : Y \to X$ be $\RR$-exponentiable,
  and let $A$ be an object in $\RR(Y)$.
  Consider the diagram
  \[\begin{tikzcd}
    A & {f^* f_* A} & {f_* A} \\
    & Y & X
    \arrow[from=1-1, to=2-2]
    \arrow["{\epsilon}"', from=1-2, to=1-1]
    \arrow["{f'}", from=1-2, to=1-3]
    \arrow[from=1-2, to=2-2]
    \arrow["\lrcorner"{anchor=center, pos=0.125}, draw=none, from=1-2, to=2-3]
    \arrow[from=1-3, to=2-3]
    \arrow["f"', from=2-2, to=2-3]
  \end{tikzcd}\]
  where $\epsilon : f^* f_* A \to A$
  is the counit of the adjunction $f^* \dashv f_*$ at $A$.
  Then there is a canonical natural isomorphism
  $f_*/A \iso {f'}_* \circ \epsilon^*$.
  \[
    \begin{tikzcd}
      {\RR(A)} & {\RR(f^* f_* A)} & {\RR(f_* A)}
      \arrow["{{\epsilon^*}}", from=1-1, to=1-2]
      \arrow["{f'_*}", from=1-2, to=1-3]
      \arrow["{f_*/A}"', from=1-1, to=1-3, bend right]
    \end{tikzcd}
  \]
\end{lemma}
\begin{proof}
  Again,
  embedding everything into presheaves $\RR(A) \to \wCC / A$,
  the result follows from \cite[Lemma 1.14]{barton2026}.
\end{proof}

Like with \Cref{lem:pullback-morphism-action}, 
it follows from \Cref{lem:distributivity} that we have a \emph{distributivity law}
or \emph{type theoretic axiom of choice}
(cf. \cite[Proposition 2.4.15]{joyal2017}).
\begin{equation}
  f_* \circ A_! \iso (f_* A)_! \circ f'_* \circ \epsilon^* 
\end{equation}
Again, in an informal type theoretic syntax, one could render it as follows.
\[
  \prod_{x : X} \sum_{y : Y(x)} A(y) \iso \sum_{f : \Pi_{X} Y} \prod_{x : X} A(f x)
\]

\begin{lemma}\label{lem:pushforward-preserves-maps}
  Let $f : Y \to X$ be $\RR$-exponentiable.
  Suppose $(\CC,\SS)$ is another preclan such that $\SS \subseteq \RR$
  and $f$ is also $\SS$-exponentiable.
  Then the pushforward functor
  \[ f_* : (\RR(Y), \SS) \to (\RR(X), \SS)\]
  is a preclan morphism.
\end{lemma}
\begin{proof}
  The proof is similar to (1) $\Rightarrow$ (4) in the proof of \Cref{theorem:exponentiable}.
  Being a right adjoint, $f_* : \RR(Y) \to \RR(X)$ preserves all pullback squares.
  It remains to show that $f_*$ preserves $\SS$-maps.
  Consider an $\SS$-map $h : A \to B$ over $Y$.
  \[\begin{tikzcd}
    A & B & Y & X
    \arrow["h", from=1-1, to=1-2]
    \arrow[from=1-2, to=1-3]
    \arrow["f", from=1-3, to=1-4]
  \end{tikzcd}\]
  Recall the diagrams from \Cref{lem:distributivity}.
  \[\begin{tikzcd}
    B & {f^* f_* B} & {f_* B} \\
    & Y & X
    \arrow[from=1-1, to=2-2]
    \arrow["{\epsilon}"', from=1-2, to=1-1]
    \arrow["{f'}", from=1-2, to=1-3]
    \arrow[from=1-2, to=2-2]
    \arrow["\lrcorner"{anchor=center, pos=0.125}, draw=none, from=1-2, to=2-3]
    \arrow[from=1-3, to=2-3]
    \arrow["f"', from=2-2, to=2-3]
  \end{tikzcd}
  \quad \quad
  \begin{tikzcd}
    {\SS(B)} & {\SS(f^* f_* B)} & {\SS(f_* B)} \\
    {\RR(B)} & {\RR(f^* f_* B)} & {\RR(f_* B)}
    \arrow["{{{\epsilon^*}}}", from=1-1, to=1-2]
    \arrow[from=1-1, to=2-1]
    \arrow["{{f'_*}}", from=1-2, to=1-3]
    \arrow[from=1-2, to=2-2]
    \arrow[from=1-3, to=2-3]
    \arrow["{{{\epsilon^*}}}", from=2-1, to=2-2]
    \arrow["{{f_*/B}}"', bend right, from=2-1, to=2-3]
    \arrow["{{f'_*}}", from=2-2, to=2-3]
  \end{tikzcd}
  \]
  Although $A$ and $B$ in $\RR(Y)$ might not be in $\SS(Y)$,
  $h : A \to B$ is an object in $\SS(B) \subseteq \RR(B)$,
  and therefore we can compute
  \[f_* h \iso f'_* \, \epsilon^* \, h\]
  as an object in $\SS(f_* B) \subseteq \RR(f_* B)$,
  making $f_* h$ an $\SS$-map.
\end{proof}

\section{Examples} \label{sec:example}

\begin{example}
  Let $\CC$ be a lex category
  and let $\RR$ consist of all morphisms in $\CC$.
  Then $(\CC,\RR)$ forms a clan.
  Functors between slices always preserve $\RR$-maps,
  making condition (4) of \Cref{theorem:exponentiable} trivial.
  Thus, $\RR$-exponentiability reduces to
  the standard notion of exponentiability of a morphism in a lex category.
  In terms of the pullback-pushforward adjunction,
  since $\RR(X) = \CC / X$ for all $X$,
  the partial right adjoint to pullback $f^* : \CC / X \to \CC / Y$ becomes
  an actual right adjoint $f_* = \Pi_f : \CC/Y \to \CC/X$.
  % Conversely,
  % any exponentiable morphism $f : Y \to X$ has a
  % (global) pullback-pushforward adjunction,
  % and $\Pi_f : \CC / Y \to \CC / X$ will preserve pullbacks --
  % since it is a right adjoint -- and will preserve $\RR$-maps trivially.
\end{example}

\begin{example}[Topological spaces, open embeddings and closed embeddings]
  \label{ex:top-opens}
  Consider the category of topological spaces $\Top$.
  We write $\OO$ for the class of open embeddings,
  i.e. those maps homeomorphic to the inclusion of an open subset,
  and $\HH$ for the class of closed embeddings,
  i.e. those maps homeomorphic to the inclusion of a closed subset.
  These form a $\tau$-preclan $(\Top,\OO,\HH)$ (\Cref{def:double-preclan}).

  It is easy to check that embeddings and open embeddings form clans in $\Top$.
  Writing $\Sub$ for the class of embeddings,
  we can define a right adjoint to $f^* : \Sub(X) \to \Sub(Y)$ as follows.
  For $U \subseteq Y$ a subspace,
  let $(\forall f)(U) \subseteq X$ denote the subspace
  \[\{\, x \in X \mid f(y) = x \Rightarrow y \in U \,\}\]
  It is easy to see that $(\forall f)(U) \in \Top/X$
  represents the functor $\Hom_{\Top/Y}(f^*(-), U)$.
  \[ \begin{tikzcd}
    {\Top / Y} & {\Sub(Y)} & {\OO(Y)} \\
    {\Top / X} & {\Sub(X)} & {\OO(X)}
    \arrow[from=1-2, to=1-1]
    \arrow[""{name=0, anchor=center, inner sep=0}, "{\forall f}", shift left=3, from=1-2, to=2-2]
    \arrow[hook', from=1-3, to=1-2]
    \arrow["{f^*}", from=2-1, to=1-1]
    \arrow[""{name=1, anchor=center, inner sep=0}, "{f^*}", shift left=3, from=2-2, to=1-2]
    \arrow[from=2-2, to=2-1]
    \arrow["{f^*}", from=2-3, to=1-3]
    \arrow[hook', from=2-3, to=2-2]
    \arrow["\dashv"{description}, draw=none, from=1, to=0]
  \end{tikzcd} \]
  
  A map $f : Y \to X$ is $\OO$-exponentiable
  if and only if it is universally closed
  (all pullbacks of $f$ are closed).
  Indeed, we show that
  the partial right adjoint condition holds at $f : Y \to X$
  if and only if $f$ is closed ($f$-image of a closed subset is closed):
  the partial right adjoint condition holds at $f$
  if and only if 
  $\forall f : \Sub(Y) \to \Sub(X)$
  sends $\OO(Y) \subseteq \Sub(Y)$ into $\OO(X) \subseteq \Sub(X)$,
  i.e. if $U \subseteq Y$ is open then 
  $(\forall f)(U) \subseteq X$ is open.
  By passage to complements, this is the same as asking that the map $f$ is closed.
  
  By a dual argument,
  the partial right adjoint condition with respect to the closed embeddings
  $\HH$ holds at $f : Y \to X$
  if and only if $f$ is an open map ($f$-image of an open subset is open).
  Since open maps are stable under pullback (unlike closed maps; see below),
  it follows that
  $f$ is $\HH$-exponentiable if and only if
  $f$ is an open map.
  \[\OO \subseteq \{\text{open maps}\} = \Exp_\HH\]
  Also, any closed embedding is universally closed.
  \[\HH \subseteq \Exp_\OO\]
  Thus, open embeddings and closed embeddings together
  form a $\tau$-preclan $(\Top,\OO,\HH)$.

  In general, a map $f : Y \to X$ is universally closed
  if and only if it is closed and the preimage of any compact subset of $X$ is compact
  \cite[Theorem 055R]{stacks-project}.
  For maps between locally compact Hausdorff spaces,
  the assumption that $f$ is closed can be omitted.
  Depending on the source, the term ``proper'' can refer to being universally closed,
  just the condition that preimages of compact subsets are compact,
  or also include the condition that $f$ is separated.

  We see that the Beck--Chevalley condition holding at $f$
  does not imply that it holds at pullbacks of $f$,
  since there are maps $f : Y \to X$ that are closed but not universally closed.
  Consider $f : \N \to 1$.
  Though $f$ is closed, the pullback $X \times \N \to X$ may not be closed.
  An $\N$-indexed union of closed subsets $U_i$ of $X$ may not be closed,
  but the disjoint union of the subsets $U_i \times \{i\}$ in $\N \times X$ is closed.
  
  We can also see that when the Beck--Chevalley condition holds at $f : Y \to X$,
  the pushforward
  \[f_* : (\OO(Y), \OO) \to (\OO(X), \OO)\]
  is automatically a preclan morphism,
  because all morphisms between objects of $\OO(X)$ are open embeddings.
  This shows that the converse of (\textit{i}) in \Cref{rmk:counterexample}
  is not true:
  $f$ satisfies Beck–-Chevalley and $f_* : (\OO(Y), \OO) \to (\OO(X), \OO)$
  is a preclan morphism, but $f$ is not $\OO$-exponentiable.
\end{example}

\begin{example}[Topological spaces and open maps]
  \label{ex:top-open-maps}
  Another choice of preclan structure $\RR$ on $\Top$
  is the class of open maps.
  In fact, this $(\Top, \RR)$ is a clan.
  We use it to give an example of
  a map $f : Y \to X$ that satisfies the Beck--Chevalley condition
  but for which $f_* : (\RR(Y), \RR) \to (\RR(X), \RR)$ is not a preclan morphism.

  Take the unique map $f : \N \to 1$,
  where $\N$ has the discrete topology.
  Since $\N$ is locally compact, $f$ is exponentiable \cite{niefield1982}.
  Moreover, every map into $1$ (or $\N$ for that matter) is an open map,
  so $\Pi_f : \Top/\N \to \Top/1 = \Top$ certainly sends $\RR(\N)$ into $\RR(1)$.
  However, we claim that $f_* = \Pi_f$ does not preserve open maps.
  Since the exponential $(-)^\N : \Top \to \Top$
  factors as $(-)^\N = \Pi_f \circ f^*$,
  and $f^* : \Top \to \Top / \N$ preserves open maps,
  it suffices to show that
  $(-)^\N = \Pi_f \circ f^*$ does not preserve open maps.
  To see this,
  take the open point $1$ in the Sierpinski space $\mathbb{S} = \{0 < 1\}$,
  which gives us an open embedding $i : \{1\} \to \mathbb{S}$ by inclusion.
  We want to show that \[i^\N : \{1\}^\N \to \mathbb{S}^\N\]
  is not an open map,
  by showing that the constant map $c_1 : \N \to \mathbb{S}$
  at $1 \in \mathbb{S}$
  is not an open point in $\mathbb{S}^\N$.

  Recall that $\mathbb{S}^\N$ is given the Scott topology \cite{niefield1982};
  in our case we can provide a basis for the Scott topology on $\mathbb{S}^\N$
  where a basic open is of the form 
  \[\{f : \Top(\N,\mathbb{S}) \st f(F) \subseteq \{1\} \}\]
  for some finite set of naturals $F \subseteq \N$.
  Suppose that the constant function $c_1$ were an open point in $\mathbb{S}^\N$.
  Then we would have
  \[ \{c_1\} = \bigcup_{i} \{f : \Top(\N,\mathbb{S}) \st f(F_i) \subseteq \{1\} \}\]
  for some finite sets $F_i \subset \N$.
  There is some $j$ such that $c_1 (F_j) \subseteq \{1\}$.
  Consider the continuous map $g : \N \to \mathbb{S}$ that is constant at $1$
  on $F_j$, and constant at $0$ elsewhere.
  Then 
  \[g \in \bigcup_{i} \{f : \Top(\N,\mathbb{S}) \st f(F_i) \subseteq \{1\} \}\]
  is in the union of opens, but $g \ne c_1$.
\end{example}

\begin{example}[Simplicial sets and Kan fibrations]
  \label{ex:sset-kan}
  Consider the category of simplicial sets $\sSet$
  with the Kan--Quillen model structure.
  Writing $\RR$ for the Kan fibrations,
  we have that $(\sSet, \RR)$ is a $\pi$-preclan \cite{joyal2017}.
  Since $\sSet$ is locally Cartesian closed,
  it is enough to check the ``Frobenius property'':
  pushforwards of fibrations along fibrations are fibrations.

  Whilst all fibrations are $\RR$-exponentiable,
  not all $\RR$-exponentiable maps are fibrations
  (an example is given below).
  The following conditions on $f : Y \to X$ are equivalent
  to $f$ being $\RR$-exponentiable.
  \begin{enumerate}
  \item
    $f^* : \sSet/X \to \sSet/Y$ preserves weak equivalences,
    i.e., in any ``two pullbacks'' diagram
    \[
      \begin{tikzcd}
        Y'' \ar[r, "t'"] \ar[d, "f''"'] \pbcorner &
        Y' \ar[r] \ar[d, "f'"'] \pbcorner & Y \ar[d, "f"] \\
        X'' \ar[r, "t"] & X' \ar[r] & X
      \end{tikzcd}
    \]
    if $t$ is a weak equivalence, then so is $t'$.
    This is Rezk's definition of a \emph{sharp map} \cite{rezk1998}.
  \item
    Every pullback of $f$ is a homotopy pullback.
  \item
    $f^* \dashv \Pi_f$ is a Quillen adjunction.
  \item
    $\Pi_f : \sSet/Y \to \sSet/X$ preserves fibrations
    ($f$ is $\RR$-exponentiable).
  \item
    $\Pi_f : \sSet/Y \to \sSet/X$ preserves fibrant objects
    (the partial right adjoint condition holds at $f$).
  \end{enumerate}

  \begin{proof}
    (1) $\Leftrightarrow$ (2) is
    proven in \cite[Proposition 2.7]{rezk1998}.

    (1) $\Leftrightarrow$ (3) $\Leftrightarrow$ (4).
    The cofibrations of $\sSet$ are the monomorphisms,
    which are stable under pullback.
    Therefore, $f^*$ always preserves cofibrations
    and by the adjunction, $\Pi_f$ always preserves trivial fibrations.
    Furthermore,
    $f^*$ preserves weak equivalences if and only if it preserves trivial cofibrations
    if and only if $\Pi_f$ preserves fibrations.

    (4) $\Rightarrow$ (5).
    Preserving fibrations implies preserving fibrant objects.

    (5) $\Rightarrow$ (1).
    A map between cofibrant objects of any model category is a weak equivalence
    just when it induces an isomorphism on
    homotopy classes of maps into every fibrant object.
    In $\sSet$ and its slices, we can use the product $A \times \Delta^1$
    as the cylinder object to compute the (left) homotopy relation.
    Let $A \in \sSet/X$, and $Z \in \sSet/Y$.
    In addition to the adjunction
    \[
      \Hom_{/Y}(f^* A, Z) \cong \Hom_{/X}(A, \Pi_f Z),
    \]
    we also have
    \[
      \Hom_{/Y}(f^*(A) \times \Delta^1, Z) =
      \Hom_{/Y}(f^*(A \times \Delta^1), Z) \cong \Hom_{/X}(A \times \Delta^1, \Pi_f Z),
    \]
    Therefore, when $Z$ is fibrant, $\Pi_f Z$ is also fibrant by (5)
    and the adjunction isomorphism passes to homotopy classes of maps
    \[
      [f^* A, Z]_{/Y} \cong [A, \Pi_f Z]_{/X}.
    \]
    This isomorphism is natural in $A$.
    So if $g : A \to B$ is a weak equivalence in $\sSet_{/X}$,
    then 
    \[g^* : [B, \Pi_f Z]_{/X} \to [A, \Pi_f Z]_{/X}\]
    is an isomorphism for all fibrant $Z \in \sSet/X$,
    and so is \[(f^*g)^* : [f^* B, Z]_{/Y} \to [f^* A, Z]_{/Y}\]
    so $f^*g : f^* A \to f^* B$ is a weak equivalence.
    This proves (1).
  \end{proof}
  
  We provide an example of an $\RR$-exponentiable map that is not a fibration.
  One can show that a weak equivalence $f : Y \to X$ is sharp
  if and only if its pullback along
  any nondegenerate simplex $\Delta^n \to X$ is again a weak equivalence,
  if and only if the pullback (object) is weakly contractible.
  The map $f : \Lambda^2_1 \to \Delta^1$
  sending vertices $0$, $1$ to $0 \in \Delta^1$ and $2$ to $1 \in \Delta^1$
  is sharp because $\Lambda^2_1$ and the fibres $(\Lambda^2_1)_0 \cong \Delta^1$,
  $(\Lambda^2_1)_1 \cong \Delta^0$ are all weakly contractible.
  However, this $f$ is not a fibration.
\end{example}

\begin{example}[Cartesian cubical sets and (unbiased) fibrations]
  \label{ex:cSet-unbiased-fibrations}
  For the following facts, we refer to \cite{awodey2026cartesian}.
  Let $\cSet$ denote the category of Cartesian cubical sets.
  Then the class of unbiased fibrations (henceforth just ``fibrations'')
  $\RR$ provides a $\pi$-preclan structure
  $(\cSet,\RR)$.
  The cubical interval $I = y([1])$ is not a fibrant,
  i.e. $I = y([1]) \to 1$ is not an $\RR$-map,
  but $I = y([1]) \to 1$ is $\RR$-exponentiable.
  It is not fibrant because the following open box in $I$ does not have a filler.
  \[ \begin{tikzcd}
    0 & 1 \\
    0 & 0
    \arrow[from=2-1, to=1-1]
    \arrow[from=2-1, to=2-2]
    \arrow[from=2-2, to=1-2]
  \end{tikzcd} \]

  To show that $I \to 1$ is $\RR$-exponentiable,
  let us first note that exponentiating by $I$ preserves fibrations.
  Let $f : A \to B$ be a fibration.
  Consider the unique map $u : 0 \to I$
  and the following ``pullback-hom'' diagram.
  \[\begin{tikzcd}
    {A^I} && \\
    & \bullet & {A^0} \\
    & {B^I} & {B^0}
    \arrow["{u \Rightarrow f}", two heads, from=1-1, to=2-2]
    \arrow["{A^u}", bend left, from=1-1, to=2-3]
    \arrow["{f^I}"', bend right, from=1-1, to=3-2]
    \arrow[from=2-2, to=2-3]
    \arrow["\iso"', from=2-2, to=3-2]
    \arrow["\lrcorner"{anchor=center, pos=0.125}, draw=none, from=2-2, to=3-3]
    \arrow["\iso", from=2-3, to=3-3]
    \arrow["{B^u}"', from=3-2, to=3-3]
  \end{tikzcd}\]
  Since $u : 0 \to I$ is a cofibration and $f : A \to B$ is a fibration,
  $u \Rightarrow f : A^I \to \bullet$ is a fibration,
  and so is $f^I : A^I \to B^I$ by the isomorphism.

  Using this fact,
  we can construct the pushforward of a fibration $f : X \to I$
  along $I \to 1$ as follows
  \[ \begin{tikzcd}
    X & {\Pi_I f} & {X^I} \\
    I & 1 & {I^I}
    \arrow["f"', two heads, from=1-1, to=2-1]
    \arrow[from=1-2, to=1-3]
    \arrow[from=1-2, to=2-2]
    \arrow["\lrcorner"{anchor=center, pos=0.125}, draw=none, from=1-2, to=2-3]
    \arrow["{f^I}", two heads, from=1-3, to=2-3]
    \arrow[from=2-1, to=2-2]
    \arrow["{\tilde{\id}}"', from=2-2, to=2-3]
  \end{tikzcd}\]
  where $\tilde{\id} : 1 \to I^I$ is the adjoint transpose 
  of the identity on $I$.
  Then $\Pi_I f \to 1$ is a pullback of a fibration,
  so is itself a fibration.

  This example will be revisited in \Cref{sec:path-types},
  to understand ``path types'' in $\cSet$.
\end{example}

% \begin{example}[Topological spaces and local homeomorphisms]
%   \label{ex:top-local-homeo}
%   Another preclan structure on topological spaces $(\Top,\RR)$ arises from
%   taking $\RR$ to be local homeomorphisms.
%   Local homeomorphisms over $X$ are equivalent to sheaves on $X$.
%   \[ \RR(X) \simeq \Sh(X) \]
%   For a continuous map $f : Y \to X$ there is always an adjunction
%   \[ \begin{tikzcd}
%     {\RR(Y)} & {\RR(X)}
%     \arrow[""{name=0, anchor=center, inner sep=0}, "{f_*}"', shift right=2, from=1-1, to=1-2]
%     \arrow[""{name=1, anchor=center, inner sep=0}, "{f^*}"', shift right=2, from=1-2, to=1-1]
%     \arrow["\dashv"{anchor=center, rotate=-90}, draw=none, from=1, to=0]
%   \end{tikzcd}\]
%   where the right adjoint $f_* : \RR(Y) \to \RR(X)$
%   is the direct image functor.
  


%   % This is closely related to the locale/($\infty$-)topos examples,
%   % via the functor sending a topological space to
%   % its underlying locale or sheaf ($\infty$-)topos.
%   % As in those examples, $\RR(X)$ is equivalent to $\Sh(X)$ for each $X$
%   % and in particular there is always a direct image functor
%   %  right adjoint to $f^* : \RR(X) \to \RR(Y)$
%   (but not necessarily satisfying the Beck--Chevalley condition).

%   In this case, we have strict inclusions
%   \[
%     \text{proper maps}
%     \subsetneq \text{$\RR$-exponentiable maps}
%     \subsetneq \text{universally closed maps} \,.
%   \]
%   As in \cite[Theorem 055R]{stacks-project}, we use ``proper map'' to mean
%   a map that is universally closed and separated.

%   \begin{itemize}
%   \item
%     $\RR$-exponentiable maps are universally closed:
%     For $X \in \Top$, the full subcategory $\OO(X) \subseteq \RR(X)$
%     of open embeddings into $X$.
%     is the subcategory of all subterminal objects of $\RR(X)$:
%     those objects $U \in \RR(X)$ whose diagonal $U \to U \times U$ is an isomorphism.
%     Note that in any preclan $(\CC, \RR)$,
%     $\RR(X) \subseteq \CC/X$ is closed under finite products
%     and therefore the condition of being subterminal
%     is the same in $\RR(X)$ and in $\CC/X$.
%     The pushforward and base change functors involved in the Beck--Chevalley condition
%     all restrict to the full subcategories $\OO(-) \subseteq \RR(-)$,
%     because they are right adjoints
%     or partial right adjoints (respectively).
%     Consequently, if $f$ satisfies the Beck--Chevalley condition with respect to $\RR(-)$,
%     then $f$ will also satisfy it with respect to $\OO(-)$.
%     By the previous example, this means that $f$ is universally closed.
%   \item
%     Proper maps are $\RR$-exponentiable:
%     This is the proper base change theorem (for sheaves of sets in topology).
%     Specifically, using $\RR(-) \simeq \Sh(-)$,
%     \cite[Theorem 055R]{stacks-project} says that
%     the Beck--Chevalley condition holds at $f$ if $f$ is proper.
%     Proper maps are stable under pullback,
%     so they are $\RR$-exponentiable.
%   \item
%     An $\RR$-exponentiable map which is not proper:
%     Let $\S$ denote the Sierpinski space.
%     The map $f : \S \to 1$ is not proper,
%     but we claim it is $\RR$-exponentiable.

%     Write $i : 1 \to \S$ for the inclusion of the closed (special) point of $\S$.
%     We will use the locally posetal structure of $\Top$:
%     for parallel maps $h_0$, $h_1 : A \to B$,
%     we write $h_0 \le h_1$ if $h_0(a)$ is a specialization of $h_1(a)$
%     (meaning $h_0(a) \in U \Rightarrow h_1(a) \in U$ for all open $U$)
%     for each $a \in A$.
%     A 2-cell $h_0 \le h_1$ determines
%     an inequality $h_0^{-1} \le h_1^{-1} : \OO(B) \to \OO(A)$
%     and so by left Kan extension
%     a natural transformation $\alpha : h_0^* \to h_1^* : \Sh(B) \to \Sh(A)$.
%     In particular, an adjunction $h \dashv k$ in $\Top$
%     gives rise to an adjunction of geometric morphisms
%     that identifies $k_*$ with $h^*$.
%     In particular, $i \dashv f$ so $f_* : \Sh(\S) \to \Sh(1)$
%     can be identified with $i^*$.
%     Likewise, the product projection $g : X \times \S \to X$
%     has a left adjoint $j = X \times i : X \to X \times \S$
%     (the product is a 2-functor, because the 2-cells of $\Top$
%     are detected by maps out of $\S$).
%     Hence for the pullback square
%     \[
%       \begin{tikzcd}
%         X \times \S \ar[r, "q"] \ar[d, "g"] & \S \ar[d, "f"] \\
%         X \ar[r, "p"] \ar[u, bend left, "j"] & 1 \ar[u, bend left, "i"]
%       \end{tikzcd}
%     \]
%     we have
%     $p^* f_* = p^* i^* = j^* q^* = g_* q^*$,
%     verifying the Beck--Chevalley condition at $f$.
%     A similar argument handles a further pullback along a map $X' \to X$.
%     We could replace $\S$ in this example
%     by any other space with a focal (= initial) point.
%   \item
%     A universally closed map which is not $\RR$-exponentiable:
%     We claim that if $f : Y \to 1$ is $\RR$-exponentiable,
%     then the global sections functor $\Gamma = f_* : \Sh(Y) \to \Sh(1) = \Set$
%     preserves filtered colimits.
%     \cite[C3.4.1]{johnstone2002} gives an example of a space $Y$
%     (two copies of $[0,1]$ glued along their common subspaces $(0,1)$)
%     which is compact, but such that $\Gamma : \Sh(Y) \to \Set$
%     does not preserve filtered colimits.

%     Moerdijk and Vermeulen prove this in the setting of 1-topoi \cite{moerdijk2000}.
%     That result cannot be applied directly here,
%     because the Beck--Chevalley condition in topoi
%     is a priori different from the one in topological spaces
%     (not every topos comes from a topological space,
%     and the products in the two categories may also differ).
%     However, one can check that the proof also goes through in the category $\Top$.
%     We omit the details for now.
%   \end{itemize}
% \end{example}


% \begin{example}[Locales and etale maps]
%   \label{ex:locale-etale}
%   The category of locales $\Loc$, like $\Top$, has several clan structures.
%   In particular, \Cref{ex:top-opens} has an analogy in this setting,
%   where the open embeddings (open regular monomorphisms) $\OO$ form a clan,
%   closed maps of locales are those that satisfy Beck-Chevalley \cite{vermeulen1994} $\OO$.
  
%   For a detailed account of proper maps of locales,
%   see \cite{vermeulen1994}.
%   We could 

%   The category of locales and \'etale maps forms a preclan $(\Loc,\EE)$.
%   Sheaves on a locale $X$ are equivalent to \'etale maps over $X$,
%   such that for a geometric morphism of locales $f : Y \to X$,
%   pullback $f^* : \EE(X) \to \EE(Y)$
%   corresponds to an inverse image functor on sheaves
%   $f^* : \Sh(X) \to \Sh(Y)$.
%   \[ \begin{tikzcd}
%     Y & {\Sh(Y)} & {\EE(Y)} & {\Loc / Y} \\
%     X & {\Sh(X)} & {\EE(X)} & {\Loc / X}
%     \arrow["f"', from=1-1, to=2-1]
%     \arrow["\simeq"{description}, draw=none, from=1-2, to=1-3]
%     \arrow[hook, from=1-3, to=1-4]
%     \arrow["{f^*}", from=2-2, to=1-2]
%     \arrow["{f^*}", from=2-3, to=1-3]
%     \arrow["\simeq"{description}, draw=none, from=2-3, to=2-2]
%     \arrow[hook, from=2-3, to=2-4]
%     \arrow["{f^*}", from=2-4, to=1-4]
%   \end{tikzcd} \]
%   The geometric morphism $f : Y \to X$
%   always produces a direct image functor $f_* : \Sh(Y) \to \Sh(X)$ on sheaves,
%   which gives us a right adjoint to pullback of \'etale maps.
%   \[ \begin{tikzcd}
%     Y & {\Sh(Y)} & {\EE(Y)} & {\Loc / Y} \\
%     X & {\Sh(X)} & {\EE(X)} & {\Loc / X}
%     \arrow["f"', from=1-1, to=2-1]
%     \arrow["\simeq"{description}, draw=none, from=1-2, to=1-3]
%     \arrow[""{name=0, anchor=center, inner sep=0}, "{f_*}", shift left=2, from=1-2, to=2-2]
%     \arrow[hook, from=1-3, to=1-4]
%     \arrow[""{name=1, anchor=center, inner sep=0}, "{f_*}", shift left=2, dashed, from=1-3, to=2-3]
%     \arrow[""{name=2, anchor=center, inner sep=0}, "{f^*}", shift left=2, from=2-2, to=1-2]
%     \arrow[""{name=3, anchor=center, inner sep=0}, "{f^*}", shift left=2, from=2-3, to=1-3]
%     \arrow["\simeq"{description}, draw=none, from=2-3, to=2-2]
%     \arrow[hook, from=2-3, to=2-4]
%     \arrow["{f^*}", from=2-4, to=1-4]
%     \arrow["\dashv"{anchor=center}, draw=none, from=2, to=0]
%     \arrow["\dashv"{anchor=center}, draw=none, from=3, to=1]
%   \end{tikzcd}\]
%   However, $f_* : \EE(Y) \to \EE(X)$ may not be a partial right adjoint to
%   $f^* : \Loc/X \to \Loc/Y$;
%   and even if it is, the same might not hold for pullbacks of $f$.
%   To see this,

%   % The condition that $f_* : \EE(Y) \to \EE(X)$ is a partial right adjoint

%   % Locales, topoi and $\infty$-topoi each form preclans
%   % whose $\RR$-maps are the etale geometric morphisms,
%   % so that $\RR(X) \simeq \Sh(X)$ for each $X$.%
%   % \footnote{
%   %   The standard notation for topoi would identify a topos $X$
%   %   and its category (logos) of sheaves $\Sh(X)$.
%   % }
%   % Write $(\Topos, \RR)$ for one of these three preclans.
%   % which is the ``inverse image'' functor usually denoted by $f^*$.
%   % The functor $f^*$ always has a right adjoint $f_* : \Sh(Y) \to \Sh(X)$,
%   % the ``direct image''.

%   A geometric morphism $f : Y \to X$ of $\infty$-topoi is called \emph{proper}
%   if the Beck--Chevalley condition holds for every pullback of $f$
%   \cite[Definition 7.3.1.4]{lurie2009}.
%   This is then equivalent to our condition that $f$ is $\RR$-exponentiable.
%   For locales and $1$-topoi, this is probably a theorem rather than the definition
%   (and the terminology may be inconsistent [?]). 
%   One justification for the terminology is the proper base change theorem (see below),
%   some version of which holds in all these settings.
% \end{example}

\begin{example}[Topological spaces and \'etale maps]
  \label{ex:topos-etale}
  ``Proper maps'' provide a prototypical example of $\RR$-exponentiability.
  They can be found in topological spaces, locales, 1-toposes,
  and $\infty$-toposes, where $\RR$ is the class of \'etale maps.
  These examples are very similar,
  and share a lot of terminology that may not be consistent across the literature;
  the word ``proper'' is used to mean various things
  depending on the context.
  For example, Lurie uses the word ``proper'' for \'etale-exponentiability
  in $\infty$-toposes \cite[Definition 7.3.1.4]{lurie2009},
  whereas Moerdijk and Vermeulen say ``tidy'' \cite[III Definition 1.2]{moerdijk2000proper}
  or ``stable BCC'' (Beck--Chevalley condition) for \'etale-exponentiability in 1-toposes,
  and reserve the word ``proper'' for a weaker condition.
  To keep things precise,
  in this example we will specifically consider \'etale maps
  (also called local homeomorphisms)
  of topological spaces.
  
  The category of topological spaces with the class of
  \'etale maps $(\Topos,\EE)$ forms a clan,
  such that \'etale maps over $X$ are
  equivalent to sheaves on $X$
  \[\EE(X) \simeq \Sh(X)\]
  for each topos $X$ \cite[Section II.6]{maclane2012}.
  For a continuous map $f : Y \to X$,
  the pullback $f^* : \Topos/X \to \Topos/Y$ preserves \'etale maps
  and so restricts to $f^* : \EE(X) \to \EE(Y)$,
  which corresponds to the ``inverse image'' functor $f^* : \Sh(X) \to \Sh(Y)$
  under the equivalence.
  There is also a ``direct image''
  functor $f_* : \Sh(Y) \to \Sh(X)$, right adjoint to
  the inverse image $f^*  : \Sh(X) \to \Sh(Y)$.
  This gives us a right adjoint to the pullback $f^* : \EE(X) \to \EE(Y)$
  of \'etale maps.
  \[ \begin{tikzcd}
    Y & {\Sh(Y)} & {\EE(Y)} & {\Top / Y} \\
    X & {\Sh(X)} & {\EE(X)} & {\Top / X}
    \arrow["f"', from=1-1, to=2-1]
    \arrow["\simeq"{description}, draw=none, from=1-2, to=1-3]
    \arrow[""{name=0, anchor=center, inner sep=0}, "{f_*}", shift left=2, from=1-2, to=2-2]
    \arrow[hook, from=1-3, to=1-4]
    \arrow[""{name=1, anchor=center, inner sep=0}, "{f_*}", shift left=2, dashed, from=1-3, to=2-3]
    \arrow[""{name=2, anchor=center, inner sep=0}, "{f^*}", shift left=2, from=2-2, to=1-2]
    \arrow[""{name=3, anchor=center, inner sep=0}, "{f^*}", shift left=2, from=2-3, to=1-3]
    \arrow["\simeq"{description}, draw=none, from=2-3, to=2-2]
    \arrow[hook, from=2-3, to=2-4]
    \arrow["{f^*}", from=2-4, to=1-4]
    \arrow["\dashv"{anchor=center}, draw=none, from=2, to=0]
    \arrow["\dashv"{anchor=center}, draw=none, from=3, to=1]
  \end{tikzcd}\]
  However, $f_* : \EE(Y) \to \EE(X)$ may not satisfy Beck--Chevalley;
  and even if it does, the same might not hold for pullbacks of $f$.
  Hence, the converse of $\textit(iii)$ in \Cref{rmk:counterexample} fails.
  Moerdijk and Vermeulen characterise those $\EE$-exponentiable maps
  as those that are ``tidy'' \cite[III, Corollary 4.9]{moerdijk2000proper}.
  %TODO example of a geometric morphism that doesn't satisfy BC?
\end{example}

% \begin{example}[Toposes and \'etale maps]
%   \label{ex:topos-etale}
%   A prototypical example of $\RR$-exponentiable maps
%   can be found in locales, (Grothendieck) 1-toposes, or $\infty$-toposes,
%   where $\RR$ are the \'etale geometric morphisms.
%   These examples are very similar,
%   and share a lot of terminology that may not be consistent across the literature.
%   In particular, the word ``proper'' is used to mean various things
%   depending on the context.
%   For example, Lurie uses the word ``proper'' for \'etale-exponentiability
%   in $\infty$-toposes \cite[Definition 7.3.1.4]{lurie2009},
%   whereas Moerdijk and Vermeulen say ``tidy'' \cite[III Definition 1.2]{moerdijk2000proper}
%   or ``stable BCC'' (Beck--Chevalley condition) for \'etale-exponentiability in 1-toposes,
%   and reserve the word ``proper'' for a weaker condition.
%   To keep things precise,
%   in this example we will consider \'etale morphisms of 1-toposes.
  
%   The category of (Grothendieck) toposes with the class of
%   \'etale geometric morphisms $(\Topos,\EE)$ forms a clan,
%   so that $\EE(X) \simeq \Sh(X)$ for each topos $X$.%
%   \footnote{
%     The standard notation for toposes would identify a topos $X$
%     and its category (logos) of sheaves $\Sh(X)$.
%   }
%   For a geometric morphism $f : Y \to X$,
%   pullback $f^* : \Topos/X \to \Topos/Y$ preserves \'etale morphisms
%   and so restricts to a functor $f^* : \EE(X) \to \EE(Y)$,
%   which corresponds to the ``inverse image'' functor $f^* : \Sh(X) \to \Sh(Y)$
%   under the equivalence.
%   The geometric morphism $f$ also comes equipped with a ``direct image''
%   functor $f_* : \Sh(Y) \to \Sh(X)$, right adjoint to
%   the inverse image $f^*  : \Sh(X) \to \Sh(Y)$.
%   This gives us a right adjoint to the pullback $f^* : \EE(X) \to \EE(Y)$
%   of \'etale maps.
%   \[ \begin{tikzcd}
%     Y & {\Sh(Y)} & {\EE(Y)} & {\Topos / Y} \\
%     X & {\Sh(X)} & {\EE(X)} & {\Topos / X}
%     \arrow["f"', from=1-1, to=2-1]
%     \arrow["\simeq"{description}, draw=none, from=1-2, to=1-3]
%     \arrow[""{name=0, anchor=center, inner sep=0}, "{f_*}", shift left=2, from=1-2, to=2-2]
%     \arrow[hook, from=1-3, to=1-4]
%     \arrow[""{name=1, anchor=center, inner sep=0}, "{f_*}", shift left=2, dashed, from=1-3, to=2-3]
%     \arrow[""{name=2, anchor=center, inner sep=0}, "{f^*}", shift left=2, from=2-2, to=1-2]
%     \arrow[""{name=3, anchor=center, inner sep=0}, "{f^*}", shift left=2, from=2-3, to=1-3]
%     \arrow["\simeq"{description}, draw=none, from=2-3, to=2-2]
%     \arrow[hook, from=2-3, to=2-4]
%     \arrow["{f^*}", from=2-4, to=1-4]
%     \arrow["\dashv"{anchor=center}, draw=none, from=2, to=0]
%     \arrow["\dashv"{anchor=center}, draw=none, from=3, to=1]
%   \end{tikzcd}\]
%   However, $f_* : \EE(Y) \to \EE(X)$ may not satisfy Beck--Chevalley;
%   and even if it does, the same might not hold for pullbacks of $f$.
%   Hence, the converse of $\textit(iii)$ in \Cref{rmk:counterexample} fails.
%   Moerdijk and Vermeulen characterise those $\EE$-exponentiable maps
%   as those that are ``tidy'' \cite[III Corollary 4.9]{moerdijk2000proper}.
% \end{example}

\begin{example}[Categories and fibrations]\label{ex:categories-and-fibrations}
  The following example will be explored in greater detail in \Cref{sec:cat-as-types}.

  The category of categories contains several preclans.
  The classes of (Grothendieck) fibrations and opfibrations are both clans.
  Opfibrations are fibration-exponentiable and vice versa
  (\Cref{prop:pushforward-in-cat}),
  together forming a $\tau$-clan (\Cref{def:double-preclan}).
  However, this example is somewhat unnatural,
  since viewing (op)fibrations as clans in $\Cat$ means that
  we are considering all functors between opfibrations,
  not just opcartesian functors (\Cref{def:morphism-split-opfibration}).

  We can restrict our attention to
  \emph{groupoidal} (op)fibrations --
  those which have groupoid fibres.
  The classes of groupoidal fibrations and
  groupoidal opfibrations form a $\tau$-preclan in $\Cat$
  (as opposed to a $\tau$-clan).
  Similarly, discrete fibrations and discrete opfibrations form another $\tau$-preclan in $\Cat$.
  % All functors between groupoidal fibrations are morphisms of fibrations,
  % so this example is somewhat more natural than the previous.

  The class of (split) isofibrations in the category of groupoids forms a $\pi$-clan.
  This will be used in \Cref{sec:hottlean} to construct a model of Martin-L\"of Type Theory,
  after \cite{hofmann1995}.
  One might expect that another model of Martin-L\"of Type Theory
  could be given by using isofibrations $\II$ or groupoidal isofibrations $\GG$ in $\Cat$.
  Unfortunately, although $(\Cat,\II)$ is a clan and $(\Cat,\GG)$ is a preclan,
  neither satisfies the $\pi$-preclan condition.
  The following (standard) example of a functor that is not
  exponentiable (with respect to all morphisms) can be used to demonstrate this fact.
  Consider the embedding $f := [0,2] : \two \to \three$,
  where $\two$ is the walking arrow,
  $\three$ is the walking composition of two arrows,
  and $f$ takes the arrow to the composition.
  % It is well known that $f$ is not exponentiable:
  % $f^* : \Cat / \three \to \Cat / \two$ does not preserve colimits (this example will demonstrate this),
  % therefore it does not have a right adjoint.
  Since $\GG$ is a preclan,
  there is a pullback functor $f^* : \GG(\three) \to \GG(\two)$.
  If $f$ were $\GG$-exponentiable, $f^*$ would have a right adjoint,
  so that $f^*$ preserves colimits.
  Consider the following pushout diagram in $\GG(\three) \subseteq \II(\three) \subseteq \Cat / \three$.
  \[
    \begin{tikzcd}
      1 & \two \\
      \two & \three
      \arrow["{[0]}", from=1-1, to=1-2]
      \arrow["{[1]}"', from=1-1, to=2-1]
      \arrow["{[1,2]}", from=1-2, to=2-2]
      \arrow["{[0,1]}"', from=2-1, to=2-2]
    \end{tikzcd}
  \]
  The pullback along $f$ produces the following square, which is not a pushout in $\GG(\two)$.
  \[
    \begin{tikzcd}
      0 & 1 \\
      1 & \two
      \arrow["{!}", from=1-1, to=1-2]
      \arrow["{!}"', from=1-1, to=2-1]
      \arrow["{[1]}", from=1-2, to=2-2]
      \arrow["{[0]}"', from=2-1, to=2-2]
    \end{tikzcd}
  \]
  Finally, note that $f$ is a groupoidal isofibration,
  hence $\GG$ is not a $\pi$-preclan.
  The same example shows that $\II$ is not a $\pi$-clan and $f$ is not an exponentiable morphism.

  It is also easy to come up with examples of exponentiable maps that are not
  exponentiable with respect to opfibrations.
  All opfibrations are exponentiable \cite{giraud1964,vidmar2018},
  so it suffices to give an opfibration that is not opfibration-exponentiable.
  Consider the (discrete) opfibrations $0 \to 1$ and $b : 1 \to \{a < b\}$.
  We can calculate the pushforward explicitly as $a : 1 \to \{a < b\}$,
  which is not an opfibration.
  \[ \begin{tikzcd}
    0 & {1 = \Pi_b 0} \\
    1 & {\{a < b\}}
    \arrow["\subseteq"', from=1-1, to=2-1]
    \arrow["a", from=1-2, to=2-2]
    \arrow["b"', from=2-1, to=2-2]
  \end{tikzcd}\]
\end{example}